\pdfmapfile{+cm.map}\pdfmapfile{+cmextra.map}\pdfmapfile{+symbols.map}
\documentclass[11pt,a4paper]{article}
\usepackage[margin=24mm]{geometry}
\usepackage{amsmath,amssymb,amsthm,mathtools,mathrsfs,booktabs,longtable,array}
\usepackage[colorlinks=true,linkcolor=blue!40!black,citecolor=blue!40!black,urlcolor=blue!40!black]{hyperref}
\usepackage{xcolor}
\hypersetup{pdftitle={Arithmetic of Gieseking Moment Lattices},pdfauthor={Keio},pdfsubject={Unified research manuscript; irrationality of V remains unresolved}}
\numberwithin{equation}{section}
\newtheorem{theorem}{Theorem}[section]
\newtheorem{lemma}[theorem]{Lemma}
\newtheorem{corollary}[theorem]{Corollary}
\newtheorem{proposition}[theorem]{Proposition}
\theoremstyle{definition}\newtheorem{definition}[theorem]{Definition}
\theoremstyle{remark}\newtheorem{remark}[theorem]{Remark}
\newcommand{\Q}{\mathbb Q}\newcommand{\Z}{\mathbb Z}\newcommand{\R}{\mathbb R}
\newcommand{\B}{\mathcal B}\newcommand{\N}{\mathcal N}\newcommand{\M}{\mathcal M}
\newcommand{\A}{\mathcal A}\newcommand{\Vd}{\mathcal V}\newcommand{\Rp}{\mathbb Z_{(p)}}
\newcommand{\calR}{\mathcal R}\newcommand{\Fp}{\mathbb F_p}\newcommand{\odd}{\operatorname{odd}}
\newcommand{\lc}{\ell_c}\newcommand{\ld}{\ell_0}\newcommand{\leone}{\ell_1}
\newcommand{\Tr}{\mathop{\rm Tay}\nolimits_{<2n,1}}
\DeclareMathOperator{\Cl}{Cl}\DeclareMathOperator{\lcm}{lcm}
\setlength{\emergencystretch}{3em}
\setlength{\parskip}{3pt}
\setcounter{tocdepth}{1}
\title{Arithmetic of Gieseking Moment Lattices:\\
Normality, Integral Saturation, and Approximation Budgets}
\author{Keio}
\date{21 October 2026}
\begin{document}
\maketitle
\begin{abstract}
We study the four formal period coordinates generated by the kernel
$(1-s^2t^2/4)^{-1/2}$ and the Gieseking volume
$V=\Cl_2(\pi/3)$. A synthesis of earlier and subsequent moment constructions
leads to several further results. Three powers of a short least common
multiple suffice for the target columns of the complete moment inverse,
settling the exponent-three question left in the supplied inverse construction.
We classify the entire first residue layer of coefficient saturation in the
no-wrap regime and close a previously undetermined two-block prime interval.
For a direct odd output space we determine the exact dyadic integral-output
index and obtain an exponential lower bound for every complete integral
output basis. We optimize the positive Bernstein approximation budget,
control exceptional primes for growing conic normalizers, and identify the
irrationality exponent encoded by a calibrated moment lattice. Finally,
prime-scale coefficient certificates exclude algebraic roots of height below
$2\log p$, giving a degree-explicit criterion for future transcendence
arguments. Complete finite-window normality, local image calculations, and
the inverse construction are included as foundations. These results resolve
specific structural and arithmetic questions within the constructions.
They do not prove irrationality or transcendence of $V$: the needed
infinite smallness estimate for genuinely integral surviving forms remains
unproved.
\end{abstract}
\vspace{4pt}
\noindent\textbf{Keywords.} Gieseking volume; Clausen function; period coordinates;
Legendre polynomials; local saturation; denominator bounds; Diophantine approximation.
\newpage
\tableofcontents
\newpage

\section{Introduction and the new conclusions}
The constant under consideration is
\begin{equation}\label{intro:V}
 V=\Cl_2(\pi/3)=\frac{3\sqrt3}{4}L(2,\chi_{-3})
 =\int_0^1\!\int_0^1\frac{ds\,dt}{\sqrt{1-s^2t^2/4}}
 =1.0149416064\ldots.
\end{equation}
The arithmetic of an explicit period construction has several separate
parts: a rational-coordinate map must be nondegenerate, its local images
must be known, its denominators must be controlled, and a nonzero integral
form must be small at the real place. Solving one of these parts does not
automatically solve the others. This paper treats their interaction in a
single moment system.

The source corpus consists of 136 earlier Gieseking notes and 47 subsequent
notes, in five named series. We use the strongest compatible results and
prove further consequences, rather than concatenate all 183 notes.
The source identifiers in the bibliography refer to supplied, unpublished
manuscripts. The word ``new'' below means a deduction made in this synthesis
relative to those supplied statements; it is not a literature-wide priority
claim. The finite checks are reproducible and supplement the proofs.

\subsection{Results established by combining the constructions}
The main advances are as follows.
\begin{enumerate}
\item \textbf{Target inverse denominator.}
For every $n$ and prime $p\ge5$, each target column of the full $4n$-column
inverse has valuation at least $-3\lfloor\log_p(2n-1)\rfloor$
(Theorem~\ref{inv3:thm:inverse-three}). The previously proved factor four
is reduced to three. This resolves the explicit unproved improvement in
Meiji R62; optimality of three is not claimed.
\item \textbf{Complete first-layer saturation.}
Theorem~\ref{sat:reiwa:classification} classifies all sources in a no-wrap
window, not just selected divided columns. Only the two boundary prime
blocks can enlarge its source span. Combining the resulting upper inclusion
with the earlier nested-image lower inclusion closes
$19/100<p/n<39/200$:
\[
 \Gamma_{\rm sat,(p)}=\Gamma_{\rm nested,(p)}
       =p^{-1}\Z_{(p)}\times p\Z_{(p)}.
\]
The polynomial exceptions and compact-interval qualification are retained
in Theorem~\ref{sat:reiwa:gapclosed}.
\item \textbf{An intrinsic integral-basis cost.}
For the direct odd rational-output basis the exact dyadic index is
\[
 2^{4n^2-2n-3\sum_{j<n}s_2(j)}.
\]
Together with the odd-prime residue ranks, this forces a geometric-mean
coefficient norm of at least
$\exp((4\log2+1/12+o(1))n)$ for every complete integral output basis
(Theorem~\ref{odd:hs:binary} and Corollary~\ref{odd:hs:height-obstruction}).
This is a restriction on the specified rational-output space.
\item \textbf{Optimized positive approximation budget.}
At full Bernstein width the best binomial order for the stated universal
denominator is $m/n=1/3$, and the sufficient control-height budget improves
from $0.9175241689\ldots$ to $1.1755429311\ldots$
(Theorem~\ref{opt:taisho:main}). The necessary infinite sequence of controls
is not constructed here.
\item \textbf{Uniform normalizer arithmetic.}
Nonzero high-index normalizers have height $\exp(O(h))$; the product of
their exceptional numerator and denominator primes through index $H$
has logarithm $O(H^2)$ (Theorem~\ref{height:keio:badprime}).
\item \textbf{The precise meaning of a calibrated lattice.}
Conditional on $V$ being irrational, its irrationality exponent satisfies
\[
 \limsup\frac{-\log\lambda_1(2a)}{\log\binom{2a}{a}}
       =1-\frac1{\mu(V)},\qquad
 \liminf\frac{-\log\lambda_2(2a)}{\log\binom{2a}{a}}
       =\frac1{\mu(V)}.
\]
The rational alternative is also determined
(Theorem~\ref{cal:thm:quant-calibration}).
\item \textbf{Algebraic specialization.}
The exact prime-scale coefficient profiles exclude every algebraic root of
height below $2\log p$, uniformly in its degree
(Theorem~\ref{alg:thm:root-height}). A resultant inequality then separates
the extra decay needed for irrationality, degree exclusion, and
transcendence (Theorem~\ref{alg:thm:degree-exclusion}).
\end{enumerate}

\subsection{External arithmetic and the remaining target}
Calegari, Dimitrov and Tang proved the $\Q$-linear independence of
$1,\zeta(2),L(2,\chi_{-3})$ \cite{CDT24}; their later work gives effective
irrationality measures for the last value \cite{CDT25}. The scalar
$3\sqrt3/4$ in \eqref{intro:V} is essential: an irrational number multiplied
by a nonzero algebraic irrational number can be rational. In particular
$V\in\Q$ would make $L(2,\chi_{-3})$ a rational multiple of $\sqrt3$,
which is compatible with the quoted theorems. No extension of their
scalar field is presumed here.

The positive-control criterion below would prove $V\notin\Q$ if its
infinite height hypothesis were met. The determinant criterion would
exclude algebraic degrees if its primitive real-place decay were met.
Neither hypothesis is proved in this paper. The unconditional results
are the explicit denominator, image, index, optimization, height, and
calibration theorems listed above.

\section{Canonical moments and conventions}
For rational polynomials $q,P$, put
\[
 \B(q,P)=\int_0^1\!\int_0^1
          \frac{q(s)P(t)}{\sqrt{1-s^2t^2/4}}\,ds\,dt.
\]
Throughout the mixed-moment sections $q$ is even. We use formal coordinates
\[
 \B(q,P)=\beta+\tau\sqrt3+\sigma\pi+\alpha V,
 \qquad (\beta,\tau,\sigma,\alpha)\in\Q^4,
\]
defined by the recurrences, rather than by an assumption of rational
independence of the four real constants. In the conic sections a different
ordering is displayed explicitly. Thus a nonzero formal vector need not
have nonzero real evaluation.

Let $J_j=\int_0^1t^j(1-t^2/4)^{-1/2}\,dt$ and
$I_j=\int_0^1-t^j\log t(1-t^2/4)^{-1/2}\,dt$. Direct integration by parts
gives
\begin{align*}
 J_0&=\pi/3,& J_1&=4-2\sqrt3,& I_0&=V,\\
 J_j&=\frac{4(j-1)}jJ_{j-2}-\frac2j\sqrt3,\\
 I_j&=\frac{4(j-1)}jI_{j-2}-\frac4jJ_{j-2}+\frac1jJ_j
 \quad(j\ge2\text{ even}).
\end{align*}
For even $i$ the mixed moment is
\[
 M_{ij}=\begin{cases}(J_j-J_i)/(i-j),&i\ne j,\\I_i,&i=j.
 \end{cases}
\]
Its target coordinate is zero off the diagonal and equals
$\binom{i}{i/2}$ on the even diagonal. These identities specify all formal
coordinates exactly. The integral identity for $V$ follows either by
integrating in $s$ and substituting $t=2\sin(u/2)$, or by the absolutely
convergent positive series
\[
 V=\sum_{r\ge0}\frac{\binom{2r}{r}}{16^r(2r+1)^2}.
\]

The $p$-adic valuation is normalized by $v_p(p)=1$ and $v_p(0)=+\infty$.
The ring $\Z_{(p)}$ is the localization of $\Z$ at $p$; completing to
$\Z_p$ does not change the finite local indices used below. A coefficient
saturation is the intersection of a rational polynomial space with the
integral coefficient lattice. All logarithms are natural. A primitive
integer polynomial has coefficient gcd one. A sufficient denominator
bound is distinguished from a least denominator.



\section{Exact finite-window normality}
\subsection{The determinant and its provenance}
The selected target is irrationality of $V=\operatorname{Cl}_2(\pi/3)$, the Gieseking volume. No such conclusion is claimed here. We use Taisho's formal moment algebra \cite{T05}. The three-block determinant of Heisei round06 \cite{H06} motivates the present extension; its branch-point multiplicity and infinity row-reduction arguments are adapted explicitly below. The new fourth block is the second endpoint sequence $d_j(0)$. Its presence changes the zero multiplicity at the origin, the factorial constant, and the resulting full-coordinate rank theorem. Heisei's later odd-column sign-regularity construction \cite{H42} is distinct and is not used.

Define
\begin{align*}
 c_{2r}&=\binom{2r}r,&c_{2r+1}&=0,\\
 d_0(x)&=0,&d_1(x)&=-2,&
 d_j(x)&=\frac{4(j-1)}j d_{j-2}(x)-\frac2j x^{j-1}\quad(j\ge2).
\end{align*}
For $n\ge1$, let $\mathcal H_n(x)$ be the $4n$ by $4n$ matrix with columns $0\le j<4n$ and four successive row blocks
\begin{equation}\label{normal:H}
 (j^k c_j)_j,\quad(j^k d_j(0))_j,\quad(j^k d_j(x))_j,\quad(j^k x^j)_j,
 \qquad 0\le k<n.
\end{equation}
Here $j^0=1$ even at $j=0$.

\begin{theorem}\label{normal:main}
For every $n\ge1$,
\begin{equation}\label{normal:formula}
 \det\mathcal H_n(x)=C_n x^{n(3n-1)}(x^2-4)^{n^2},
\end{equation}
where the positive rational constant is
\begin{equation}\label{normal:constant}
 C_n=2^{4n^2-2n}
 \left(\prod_{i=0}^{n-1}\frac{(i!)^2((2i)!)^2}{(2n+2i+1)!}\right)
 \prod_{j=1}^{2n-1}j!.
\end{equation}
Thus $\mathcal H_n(x)$ is nonsingular exactly when $x\notin\{0,2,-2\}$.
\end{theorem}
We prove the formula by three multiplicities, a degree bound, and an explicitly nonzero leading coefficient. All local series manipulations are formal over $\Q$.

\subsection{The two branch points and infinity}
Induction in the recurrence gives
\begin{equation}\label{normal:ode}
 (4-x^2)d_j'(x)-x d_j(x)=2(x^j-c_j),\qquad
 d_{2r+1}(0)=-\frac{2\,16^r(r!)^2}{(2r+1)!}.
\end{equation}
Fix $b\in\{2,-2\}$ and write $z=x-b$. There is a unique formal series $a(x)$ at $b$ satisfying
$(4-x^2)a'-xa=-2$: the coefficient determining $[z^r]a$ is $-b(2r+1)\ne0$.
Subtract $a$ times the first block from the third and divide column $j$ by $x^j$, a unit locally. Put $F_j=x^{-j}(d_j-a c_j)$. Then
\[
 (4-x^2)F_j'+\left(j(4-x^2)/x-x\right)F_j=2.
\]
The coefficient of $z^r$ in $F_j$ is a polynomial in $j$ of degree at most $r$, by induction; every occurrence of $j$ in its recursion is multiplied by a series vanishing at $b$.
The fourth block is now $v_k=(j^k)_j$, $k<n$. After quotienting by these rows, a nonzero determinant term from the third block needs $n$ distinct degrees at least $n$. Its total order is therefore at least
\[
 \sum_{s=n}^{2n-1}s-\sum_{k=0}^{n-1}k=n^2.
\]
The other two blocks do not affect this vanishing argument. Hence both $(x-2)^{n^2}$ and $(x+2)^{n^2}$ divide the determinant.

For the degree bound, divide column $j$ by $x^j$ at infinity. In the first and third blocks replace $j^k$ by
$P_k(j)=\prod_{s=0}^{k-1}(j-2s)$; in the second use
$Q_k(j)=\prod_{s=0}^{k-1}(j-2s-1)$. These monic triangular basis changes have determinant one. The first and second block rows then have orders $O(x^{-2k})$ and $O(x^{-2k-1})$ respectively, by their parity and the displayed roots.

As in \cite{H06}, the third row can be reduced to order $O(x^{-2k-1})$ by subtracting multiples of first- and fourth-block rows. Here are the details needed to avoid a degree assumption. Iteration of the recurrence gives
\[
 [x^{j-2r-1}]d_j(x)
 =a_r(j):=-2\,4^r
   \frac{\prod_{s=0}^{r-1}(j-2s-1)}{\prod_{s=0}^{r}(j-2s)}
 \quad(j\ge2r+1).
\]
For $r<k$, $P_k(z)a_r(z)$ is a polynomial of degree $k-1$. Subtract its row, multiplied by $x^{-2r-1}$, using the fourth block. The only exceptional indices are even $j=0,2,\ldots,2k-2$, where the rational expression originally had a removable product of a zero and a pole. The resulting exceptional values are removed exactly with the first $k$ rows of the first block: their matrix on those columns is triangular with diagonal $P_s(2s)c_{2s}x^{-2s}\ne0$. Successive solution uses multipliers $O(x^{-1})$. On all other even columns these corrections are $O(x^{-2k-1})$, and on odd columns they vanish. This proves the claimed row order.

The divided determinant has order at most
$x^{-\sum_{k<n}(2k+2k+1+2k+1)}=x^{-3n^2+n}$.
Restoring the columns gives
\begin{equation}\label{normal:degree}
 \deg\det\mathcal H_n\le\sum_{j=0}^{4n-1}j-3n^2+n=5n^2-n.
\end{equation}

\subsection{The origin and the factorial constant}
This is the additional step specific to the four-block problem. Solving \eqref{normal:ode} formally at zero gives
\[
 d_j(x)=b(x)d_j(0)+a(x)c_j+
 b(x)\int_0^x\frac{t^j}{\sqrt{4-t^2}}\,dt,
\]
where
\[
 b(x)=\frac2{\sqrt{4-x^2}},\qquad
 a(x)=-\frac2{\sqrt{4-x^2}}\int_0^x\frac{dt}{\sqrt{4-t^2}}.
\]
Subtract the first two terms from the third block. Put $w_r=\binom{2r}r/16^r$. The remaining entry has expansion
\begin{equation}\label{normal:origin}
 x^{j+1}\sum_{r\ge0}R_r(j)x^{2r},\qquad
 R_r(z)=\frac12\sum_{v=0}^r\frac{w_{r-v}w_v}{z+2v+1}.
\end{equation}
Replace the $k$th polynomial row factor in this block by
$S_k(z)=\prod_{v=0}^{k-1}(z+2v+1)$.
For $r<k$, $S_kR_r$ is a polynomial of degree at most $k-1$, so its term can be subtracted using the fourth block. At $r=k$, subtract the polynomial part as well. The leading remaining row is
\begin{equation}\label{normal:leadingrow}
 x^{j+2k+1}\frac{\lambda_k}{j+2k+1},\qquad
 \lambda_k=\tfrac12 w_k S_k(-2k-1)
 =\frac{(-1)^k(2k)!}{2\,8^k k!}.
\end{equation}
All later terms have at least two additional powers of $x$.

Expand along the last $2n$ rows. For their column set $J$, their common factors contribute at least
\[
 \sum_{j\in J}j+\sum_{k=0}^{n-1}(2k+1)
 \ge n(2n-1)+n^2=n(3n-1).
\]
Equality forces $J=\{0,\ldots,2n-1\}$. The leading determinant on these columns is $\prod_{k<n}\lambda_k$ times
\[
 \mathcal C_n=\det\begin{pmatrix}
 ((j+2k+1)^{-1})_{0\le k<n,\,0\le j<2n}\\
 (j^k)_{0\le k<n,\,0\le j<2n}
 \end{pmatrix}
 =\frac{\Vd(0,\ldots,2n-1)\Vd(1,3,\ldots,2n-1)}
 {\prod_{j=0}^{2n-1}\prod_{k=0}^{n-1}(j+2k+1)}>0.
\]
Here $\Vd$ is the increasing-node Vandermonde product. The identity follows by multiplying the columns by $\prod_{k<n}(j+2k+1)$ and comparing the resulting polynomial evaluation determinant. Alternatively, nonvanishing alone follows because a rational function $P(z)+\sum_{k<n}a_k/(z+2k+1)$ with $\deg P<n$ cannot have $2n$ distinct nonnegative zeros unless it is zero.

The remaining columns split into
$E=(2n,2n+2,\ldots,4n-2)$ and $O=(2n+1,2n+3,\ldots,4n-1)$.
Their first two block minors contribute
$\Vd(E)\Vd(O)\prod_{j\in E}c_j\prod_{j\in O}d_j(0)$, up to the interlacing sign $(-1)^{n(n-1)/2}$. The sign from $\prod\lambda_k$ is the same, so the leading coefficient at zero has sign $(-1)^n$ and is nonzero. Dividing it by $(-4)^{n^2}$ yields a positive $C_n$. Substitution of \eqref{normal:ode}, \eqref{normal:leadingrow} and the Vandermonde products simplifies exactly to \eqref{normal:constant}.

We have found coprime factors with total degree
$n(3n-1)+2n^2=5n^2-n$, equal to \eqref{normal:degree}, and a nonzero coefficient at that order. This proves Theorem~\ref{normal:main}.

\subsection{Minimal-window formal moment normality}
For clarity, the actual bilinear integral and its one-dimensional moments are
\[
 \B(F,G)=\int_0^1\!\int_0^1\frac{F(s)G(t)}{\sqrt{1-s^2t^2/4}}\,ds\,dt,
 \quad J_j=\int_0^1\frac{t^j\,dt}{\sqrt{1-t^2/4}},
 \quad I_j=\int_0^1\frac{-t^j\log t\,dt}{\sqrt{1-t^2/4}}.
\]
Here $I_0=V$, and integration by parts gives
\[
 I_j=\frac{4(j-1)}j I_{j-2}-\frac4jJ_{j-2}+\frac1jJ_j
 \qquad(j\ge2\text{ even}).
\]
For even $i$, the entries $M_{ij}=\B(t^i,t^j)$ have the moment algebra
\[
 J_j=\frac\pi3c_j+\sqrt3\,d_j(1)-2d_j(0),\qquad
 M_{ij}=\begin{cases}(J_j-J_i)/(i-j),&i\ne j,\\I_i,&i=j,\end{cases}
 \qquad [V]I_i=c_i.
\]
It is generated by the recurrences in \cite{T05}; it is a formal coordinate convention, not an independence claim. For an increasing set $A=\{a_0,\ldots,a_{n-1}\}$ of even indices in $\{0,\ldots,4n-1\}$, write the coordinate matrices on these rows and all $4n$ columns as $R,T,P,C$. Thus the specified moment entries are $R+\sqrt3 T+\pi P+VC$.

\begin{theorem}\label{normal:normal}
For every such $A$, the square matrix
$\mathcal M_A=(R^{\mathsf T},T^{\mathsf T},P^{\mathsf T},C^{\mathsf T})^{\mathsf T}$ is invertible. More precisely,
\begin{equation}\label{normal:arbitrary}
 |\det\mathcal M_A|
 =2^n3^{n^2-n}C_n\,
 \frac{\Vd(A)^4\prod_{a\in A}c_a}
      {\prod_{a\in A}a!(4n-1-a)!}.
\end{equation}
For $A=(0,2,\ldots,2n-2)$, its signed determinant is
\begin{equation}\label{normal:momentdet}
 \det\mathcal M_A=(-1)^n2^{6n^2-3n}3^{n^2-n}
 \left(\prod_{i=0}^{n-1}\frac{(i!)^4((2i)!)^2}{((2n+2i+1)!)^2}\right)
 \prod_{j=1}^{2n-1}j!.
\end{equation}
In this case
\begin{equation}\label{normal:rate}
 \log|\det\mathcal M_A|=n^2\log(3/16)+O(n\log(n+1)).
\end{equation}
\end{theorem}
\begin{proof}
Let $D(z)=\prod_{a\in A}(z-a)$ and $O=\{0,\ldots,4n-1\}\setminus A$. If $\mathcal M_Aq=0$, then $Cq=0$ forces $q_a=0$ on $A$. Construct
\[
 F(t)=\sum_{j\in O}\frac{q_j}{D(j)}t^j+\sum_{a\in A}u_at^a
\]
with $\theta^kF(1)=0$ for $k<n$, where $\theta=t\,d/dt$. The $u_a$ exist uniquely by the Vandermonde determinant. For $L_a(z)=D(z)/(z-a)$,
\[
 L_a(\theta)F(t)=\sum_{j\in O}\frac{q_j}{j-a}(t^j-t^a).
\]
Applying each coefficient functional $t^j\mapsto c_j,d_j(0),d_j(1)$ gives zero by $Pq=Rq=Tq=0$. The $L_a$ form a basis of polynomials of degree below $n$. Thus $F$ is annihilated by all four blocks of $\mathcal H_n(1)$, which is nonsingular, so $q=0$.

For completeness, keep determinants in this same change of variables. Let $\mathcal A$ be the $3n$ by $3n$ divided-difference matrix on columns $O$ with row blocks
\[
 \frac{c_j-c_a}{a-j},\qquad
 \frac{d_j(0)-d_a(0)}{a-j},\qquad
 \frac{d_j(1)-d_a(1)}{a-j}.
\]
Changing the four polynomial row bases from monomials to $L_a$ has determinant $\Vd(A)^4$. Changing the coefficient variables of $F$ to $(q_j)_{j\in O}$ and $(L_a(\theta)F)(1)$ has absolute determinant $\Vd(A)^2|\prod_{j\in O}D(j)|$. Eliminating the last variables therefore gives
\[
 |\det\mathcal A|=|\det\mathcal H_n(1)|\,
 \Vd(A)^2\Big/\left|\prod_{j\in O}D(j)\right|.
\]
The coordinate scalings are $-2,1/3,1$ on the $d(0),c,d(1)$ blocks, and the removed diagonal block contributes $\prod c_a$.
Finally
\[
 \left|\prod_{j\in O}D(j)\right|
 =\frac{\prod_{a\in A}a!(4n-1-a)!}{\Vd(A)^2}.
\]
This proves \eqref{normal:arbitrary}. For the stated even consecutive set, expansion along the last block contributes the sign $(-1)^{n(9n-3)/2}$. The coordinate scalings, the sign of $\det\mathcal H_n(1)$, and the product $\prod_{j\in O}D(j)$ contribute respectively $(-1)^n$, $(-1)^n$, and $(-1)^{n(n-1)/2}$. Their product is $(-1)^n$; substituting \eqref{normal:constant} gives \eqref{normal:momentdet}.
Stirling's estimate $\log(m!)=m\log m-m+O(\log(m+1))$, summed over the displayed products, cancels all $n^2\log n$ terms. The remaining coefficient is $\log3-4\log2$, proving \eqref{normal:rate}.
\end{proof}


\subsection{The exact rational normalization}
The small real size of the determinant and the size of its reduced numerator are very different. The following separates them precisely.
\begin{theorem}\label{normal:arithmetic}
Define
\[
 Q_n=\frac{\prod_{i=0}^{n-1}((2n+2i+1)!)^2}
 {\left(\prod_{i=0}^{n-1}(i!)^4((2i)!)^2\right)
   \left(\prod_{j=1}^{2n-1}j!\right)}.
\]
Then $Q_n$ is an integer. If the determinant in \eqref{normal:momentdet} is written as $(-1)^n a_n/b_n$ in lowest terms, with $a_n,b_n>0$, then $a_n$ has no prime factor other than $2$ and $3$, and
\begin{align*}
 \log a_n&=(6\log2+\log3)n^2+O(n\log(n+1)),\\
 \log b_n&=10\log2\,n^2+O(n\log(n+1)).
\end{align*}
For each fixed prime $p$, one also has $v_p(Q_n)=O_p(n\log(n+1))$.
\end{theorem}
\begin{proof}
For integers $q\ge2$, put
\[
 F_n(q)=4\sum_{i<n}\left\lfloor\frac iq\right\rfloor
 +2\sum_{i<n}\left\lfloor\frac{2i}q\right\rfloor
 +\sum_{j<2n}\left\lfloor\frac jq\right\rfloor
 -2\sum_{i<n}\left\lfloor\frac{2n+2i+1}q\right\rfloor.
\]
We prove $F_n(q)\le0$. With $g=\gcd(2,q)$, summing the four arithmetic floor progressions gives the exact shift identity
\begin{equation}\label{normal:floorshift}
 F_{n+q}(q)-F_n(q)=-3q+2g-4<0.
\end{equation}
For explicit verification, the increments of the four unweighted sums in their displayed order are
\[
 n,\qquad 2n+\frac{q+g-2}{2},\qquad
 4n+q,\qquad 6n+\frac{5q-g+2}{2}.
\]
Thus it suffices to consider $0\le r<q$. Set
\[
 T_r=\sum_{k=1}^3\min\bigl(r,\max(0,2r-\lfloor kq/2\rfloor)\bigr).
\]
Counting the indices for which $2r+2i+1\ge kq$ gives
\[
 F_r(q)=2\max(0,r-\lceil q/2\rceil)+\max(0,2r-q)-2T_r.
\]
For $r\le q/2$ the two positive terms vanish. For $r>q/2$, the first term in $T_r$ equals $r$. The sum of the two positive terms is $4r-2q$ when $q$ is even and $4r-2q-1$ when $q$ is odd. Consequently $F_r(q)<0$ in the latter range as well. This proves the claim for every $n,q$.

Legendre's factorial valuation formula now gives
\[
 v_p(Q_n)=-\sum_{k\ge1}F_n(p^k)\ge0
\]
for every prime $p$, proving integrality. These same expressions give a quantitative bound: for $r<q$ one has $|F_r(q)|\le6r$, and \eqref{normal:floorshift} shows $|F_n(q)|\le11n$ for every $n\ge1,q\ge2$. Also $F_n(q)=0$ for $q>4n-1$. Therefore $v_p(Q_n)=O_p(n\log(n+1))$.

The determinant is exactly
\[
 (-1)^n\frac{2^{6n^2-3n}3^{n^2-n}}{Q_n}.
\]
All numerator primes other than $2,3$ are excluded. The fixed-prime valuation bounds show
$v_2(a_n)=6n^2+O(n\log(n+1))$ and
$v_3(a_n)=n^2+O(n\log(n+1))$.
This proves the numerator estimate. Subtracting \eqref{normal:rate} gives the denominator estimate. No assertion about an individual invariant factor or a target-specialized determinant follows from these product valuations.
\end{proof}

\subsection{Consequences and exact scope}
Four coordinate prescriptions for each of $n$ rows require at least $4n$ coefficient variables. Theorem~\ref{normal:normal} attains that minimal column count. In particular, it gives unique rational polynomials of degree below $4n$ realizing any prescribed four-coordinate output. Rational target-free dual columns with actual integral outputs $\delta_{ij}$ exist in this window. Neither their signs nor small integer normalizations follow.

If $W(s)\ne0$ is an even rational polynomial of degree $2m$, consider the rows $s^{2i}W(s)$, $0\le i<n$. They have full four-coordinate row rank $4n$ on columns of degrees $0,\ldots,4(n+m)-1$. Indeed their coefficient matrix inside the $n+m$ ordinary even rows has rank $n$; apply its four diagonal copies to the invertible matrix of Theorem~\ref{normal:normal}. No positivity or root-location hypothesis on $W$ is needed. This does \emph{not} settle nuisance normality in the smaller original $3n$-column weighted window.

Formula \eqref{normal:momentdet} also shows that the complete coordinate map is invertible over $\mathbb Z_{(p)}$ for every prime $p>4n$: its entries are integral and its determinant is a unit. It does not determine its Smith factors at smaller primes. A small real determinant of a rational coordinate matrix is not an irrationality criterion; denominators, target specialization and lattice shape are separate issues.




\section{Constructing the complete moment inverse}
\subsection{The map and the theorem}
Let $V=\operatorname{Cl}_2(\pi/3)$. Define formal moments
\begin{align*}
 J_0&=\pi/3,&J_1&=4-2\sqrt3,&I_0&=V,\\
 J_j&=4(j-1)J_{j-2}/j-2\sqrt3/j,\\
 I_i&=4(i-1)I_{i-2}/i-4J_{i-2}/i+J_i/i\quad(i\ge2\text{ even}).
\end{align*}
For even $i$, set $M_{ij}=(J_j-J_i)/(i-j)$ off the diagonal and $M_{ii}=I_i$.
The matrix $\M_n$ has columns $0\le j<4n$, rows $i=0,2,\ldots,2n-2$, and coordinate blocks $(\beta,\tau,\sigma,\alpha)$ in the formal basis $(1,\sqrt3,\pi,V)$. Formal coordinates do not assume arithmetic independence of these numbers. These moments represent the double integrals with integrand $s^it^j/\sqrt{1-s^2t^2/4}$ on the unit square.

Write $\Omega_m$ for the odd part of $\operatorname{lcm}(1,\ldots,m)$, with $\Omega_0=1$, and put $L_p=\lfloor\log_p(2n-1)\rfloor$ and $L_p^0=\lfloor\log_p(2n-2)\rfloor$, the latter defined as $0$ for $n=1$.

\begin{theorem}\label{inverse:main}
The matrix $\M_n$ has the explicit inverse constructed below. For a column $z$ of that inverse:
\begin{enumerate}
\item if its input is a unit $\beta$ coordinate, its common odd denominator divides $\Omega_{2n-2}$;
\item for a unit $\tau$ coordinate and $p\ge5$, every entry has valuation at least $-L_p^0$;
\item for a unit $\sigma$ coordinate and $p\ge5$, every entry has valuation at least $-2L_p$;
\item for a unit $\alpha$ coordinate and $p\ge5$, every entry has valuation at least $-4L_p$.
\end{enumerate}
Consequently the common inverse denominator away from $2$ and $3$ divides $\operatorname{lcm}(1,\ldots,2n-1)^4$.
\end{theorem}
The target remains unresolved. Round61 proved exponential-height rational lifts and the weaker all-column bound $-8L_p$ \cite{M61}. The new ingredient is a primitive-level cancellation, not a new lattice-minimum bound. Heisei round 48v2 concerns a different odd-only output map \cite{H48}; its sharp binary denominator and removed prime band are not transferred here.

\subsection{A self-contained endpoint reduction}
Put $\theta=t\partial_t$, $c_{2r}=\binom{2r}r$, $c_{2r+1}=0$, and
\[
 d_0(x)=0,\quad d_1(x)=-2,\qquad
 d_j(x)=4(j-1)d_{j-2}(x)/j-2x^{j-1}/j.
\]
For $F=\sum f_jt^j$ set $\lc(F)=\sum f_jc_j$, $\ld(F)=\sum f_jd_j(0)$ and $\leone(F)=\sum f_jd_j(1)$.
Direct induction gives
\begin{gather}
 J_j=\pi c_j/3+\sqrt3d_j(1)-2d_j(0),\qquad
 (4-x^2)d_j'(x)-xd_j(x)=2(x^j-c_j),\label{inverse:differential}\\
 d_{2r+1}(0)=-2\,16^r(r!)^2/(2r+1)!.
\end{gather}
Let $D_n(z)=\prod_{h=0}^{n-1}(z-2h)$ and $L_i(z)=D_n(z)/(z-2i)$. If $X\in\Q[t]_{<4n}$ vanishes to order $n$ at $1$, then $q=D_n(\theta)X/n!$ has $\alpha=0$ and
\begin{equation}\label{inverse:coordinates}
 \beta_i(q)=2\ld(L_i(\theta)X)/n!,\quad
 \sigma_i(q)=-\lc(L_i(\theta)X)/(3n!),\quad
 \tau_i(q)=-\leone(L_i(\theta)X)/n!.
\end{equation}
Indeed $q_{2i}=0$, and $L_i(\theta)X(1)=0$ eliminates the subtracted $J_{2i}$ in the divided difference.

For $\deg R<2n$, use the parity Rodrigues operator
\begin{equation}\label{inverse:rod}
 (\A_nR)(t)=\frac{t(4-t^2)^{n-1/2}}{n!}
 \left.\partial_x^n\big((4-x)^{1/2}x^{n-1/2}R(\sqrt x)\big)\right|_{x=t^2}.
\end{equation}
It is polynomial and dyadic rational, since
\begin{equation}\label{inverse:dyadic}
 \A_n(t^r)=\sum_{k=0}^n(-1)^k\binom{1/2}{k}
 \binom{n-1/2+r/2}{n-k}t^{r+2k}(4-t^2)^{n-k}.
\end{equation}
Its image is exactly the joint kernel of $\lc(\theta^kF)$ and $\ld(\theta^kF)$, $0\le k<n$. To verify this, multiply its coefficient at $t^{r+2j}$ by the relevant parity weight $c_{r+2j}$ or $d_{r+2j}(0)$. The result is a fixed constant times $(-1)^j\binom nj$, so every degree-below-$n$ moment vanishes. The leading coefficients $(-1)^n\binom{n+r/2}{n}$ are nonzero; dimension counting, using two ordinary parity Vandermonde matrices, gives equality.

For prescribed $\beta=b$, $\sigma=s$, choose $F\in\Q[t]_{<2n}$ by
\begin{align}
 F_{2i}&=-3n!s_i/(c_{2i}D_n'(2i)),\label{inverse:particular}\\
 F_{2j+1}&=\frac{n!}{2d_{2j+1}(0)}
 \sum_{i=0}^{n-1}\frac{b_i\Lambda_j(2i)}{D_n'(2i)},\qquad
 \Lambda_j(z)=\prod_{h\ne j}\frac{z-(2h+1)}{2(j-h)}.\nonumber
\end{align}
Lagrange interpolation gives $\lc(L_iF)=-3n!s_i$ and $\ld(L_iF)=n!b_i/2$. Here and below $L_iF$ means $L_i(\theta)F$.

We need one endpoint identity. Write $N=n-1$, $a_i=(-2)^Ni!(N-i)!$, and
\[
 v_i(x)=\sum_{h=i}^{N}\binom{2h}h\binom hi(4-x)^{N-h}.
\]
The adjoint identity, attributed to Heisei round43 \cite{H43}, is
\begin{equation}\label{inverse:adjoint}
 L_i(-1-\theta)(4-t^2)^{-1/2}
   =a_i v_i(t^2)/(4-t^2)^{n-1/2}.
\end{equation}
For completeness, expand $L_i$ in $D_k(z)$ and apply
$D_k(-1-2x\partial_x)(4-x)^{-1/2}
 =(-2)^k(2k)!/[k!(4-x)^{k+1/2}]$; comparison gives the displayed $v_i$.
Let $B_i(F)$ denote twice the boundary-at $1$ term in integration by parts of
$\int_0^1 L_iF/\sqrt{4-t^2}\,dt$, divided by $\sqrt3$.
It is rational, depends only on the first $N$ jets of $F$ at $1$, and is0 when those jets vanish. There is no boundary term at $0$: in the expansion in $t^k\partial_t^k$, the required derivatives of $t^k/\sqrt{4-t^2}$ of orders below $k$ vanish there.
Integrating \eqref{inverse:differential} and then using \eqref{inverse:adjoint} yields
\begin{equation}\label{inverse:endpoint}
 \frac{\leone(L_iF)-B_i(F)}{a_i}
 =\frac{2\ld(L_iF)-\pi\lc(L_iF)/3}{\sqrt3a_i}
  +\frac1{n!}\int_0^1 v_i(x)U^{(n)}(x)\,dx,
\end{equation}
whenever
\begin{equation}\label{inverse:primitive}
 U^{(n)}(x)=\frac{n!F(\sqrt x)}{\sqrt3\sqrt x(4-x)^{n-1/2}}.
\end{equation}

\begin{lemma}\label{inverse:jetlemma}
Suppose $U$ satisfies \eqref{inverse:primitive}, has rational jets at $1$, and has finite derivatives through order $n-1$ at $0$ lying in $(\Q+\pi\Q)/\sqrt3$. Set
\[
 R=-\Tr\left(t^{1-2n}\sqrt{\frac3{4-t^2}}\,U(t^2)\right),\quad
 X=F+\A_nR,
\]
where $\Tr$ is the Taylor polynomial at $1$ of degree below $2n$.
Then $X$ vanishes to order $n$ at $1$, and $\leone(L_iX)=0$ for every $i<n$.
\end{lemma}
\begin{proof}
Put $U_R(x)=x^{n-1/2}\sqrt{(4-x)/3}\,R(\sqrt x)$.
The first $2n$ jets at $1$ of $U+U_R$ vanish. Equation \eqref{inverse:rod} therefore makes the first $n$ jets of $X$ vanish. Integrating the last term of \eqref{inverse:endpoint} by parts $n$ times is legitimate: $v_i$ has degree $n-1$, and the prescribed origin traces are finite. Its endpoint-at $1$ contribution is rational. All origin terms and the first term on the right are in $(\Q+\pi\Q)/\sqrt3$. The left side is rational. Since $1,\sqrt3,\pi$ are linearly independent over $\Q$, those nonrational terms cancel and the left side equals just the endpoint-at $1$ contribution. Apply this to $U+U_R$ and $X$. The endpoint jets and $B_i(X)$ vanish, proving the assertion. The traces of $U_R$ at $0$ through order $n-1$ are $0$.
\end{proof}

\subsection{Odd Laurent preimages}
For $F=t^{2j+1}$, $0\le j<n$, take
\begin{equation}\label{inverse:oddprimitive}
 U(x)=\frac1{\sqrt3}\sum_{h=0}^j
 \frac{(-1)^{n+h}\binom jh4^{j-h}(4-x)^{h+1/2}}
      {\binom{h+1/2}{n}}.
\end{equation}
Differentiating $n$ times proves \eqref{inverse:primitive}. Its jets at $1$ are rational; its origin traces are rational multiples of $1/\sqrt3$. More concretely, the expression whose Taylor polynomial is used in Lemma~\ref{inverse:jetlemma} is the Laurent polynomial
\begin{equation}\label{inverse:Laurent}
 t^{1-2n}\sum_{h=0}^j
 \frac{(-1)^{n+h}\binom jh4^{j-h}(4-t^2)^h}
      {\binom{h+1/2}{n}}.
\end{equation}
This gives an especially short exact construction for all $\beta$ inverse columns.

\subsection{Even centered-angle preimages}
For $m\le2n-2$, let $H_{n,m}(u)$ be the primitive, vanishing at $0$, of
$u^{2m}(1+u^2)^{n-m-2}$ obtained by the following finite rules:
\[
 u^{2m}(1+u^2)^{n-m-2}
 =\sum_{j=0}^{m}(-1)^j\binom mj(1+u^2)^{n-2-j}.
\]
For nonnegative exponents expand binomially and integrate each even monomial. For exponent $-k<0$, use
\begin{align}
 \int\frac{du}{(1+u^2)^k}
 &=A_k\arctan u+
 A_k\sum_{r=1}^{k-1}\frac{4^r}{2r\binom{2r}r}
                \frac{u}{(1+u^2)^r},\label{inverse:primitiveR}\\
 A_k&=\binom{2k-2}{k-1}/4^{k-1}.\nonumber
\end{align}
This follows by one integration-by-parts recurrence, beginning with $k=1$.
Write $H_{n,m}^{\circ}$ for the same expression with every $\arctan u$ replaced by $\arctan u-\pi/6$.
For $F=t^{2j}$, $j<n$, define
\begin{equation}\label{inverse:evenprimitive}
 U(x)=\frac{2n}{\sqrt3}\sum_{h=0}^{n-1}(-1)^h\binom{n-1}h
       4^{j+h-n+1}x^{n-1-h}
 H_{n,j+h}^{\circ}\left(\sqrt{\frac{x}{4-x}}\right).
\end{equation}
Without the circle this is exactly
\[
 \frac n{\sqrt3}\int_0^x(x-z)^{n-1}
                   z^{j-1/2}(4-z)^{1/2-n}\,dz.
\]
Indeed substitute $z=4u^2/(1+u^2)$ after expanding $(x-z)^{n-1}$. Thus it satisfies \eqref{inverse:primitive} and has zero origin traces through order $n-1$.
Centering subtracts only a polynomial of degree below $n$, times $\pi/\sqrt3$; the differential equation is unchanged. Its origin traces therefore have the form required in Lemma~\ref{inverse:jetlemma}. At $x=1$ the centered angle vanishes, all algebraic terms normalized by $1/\sqrt3$ are rational, and all higher Taylor coefficients are rational. This completes the inverse construction with $\tau=0$ for arbitrary $\beta,\sigma$ inputs.

\subsection{The remaining output blocks}
For pure $\tau=u$, put $F=0$ and prescribe a polynomial $U$ of degree below $n$ by
\begin{align*}
 y_i&=-n!u_i/a_i,&
 z_h&=\frac1{\binom{2h}h}\sum_{i=h}^{N}(-1)^{i-h}\binom ih y_i,\\
 [y^k]U(1+y)&=\sum_{h=k}^{N}\binom hk(-3)^{h-k}
                                  n\binom Nh z_h.
\end{align*}
Set $R=\Tr(t^{1-2n}\sqrt{3/(4-t^2)}U(t^2))$ and $q=D_n(\theta)\A_nR/n!$; here the sign is positive. Formula \eqref{inverse:endpoint} with zero origin traces shows the desired $\tau$ outputs. Explicitly the endpoint functional of $(4-x)^{N-h}$ is
$\sum_{k=h}^{N}\binom kh3^{k-h}[y^k]U(1+y)/(n\binom Nh)$; the displayed two Pascal inversions solve it. Since $U^{(n)}=0$ and the first $2n$ jets are retained, $\A_nR$ has the required order-$n$ zero at $1$. The other output blocks vanish.

Finally, for a unit $\alpha_i$ output start with $q_0=t^{2i}/c_{2i}$. Its $\alpha$ outputs are the desired unit vector. Its $\beta$ outputs are $0$. Subtract its $\tau,\sigma$ outputs using the already constructed columns. This proves invertibility without any unproved period relation or numerical rank test.

\subsection{Arithmetic of the primitive construction}
We record the elementary counting facts used below. For an odd prime $p$ and $0\le h<n$,
\begin{equation}\label{inverse:counts}
 v_p\binom{h+1/2}{n}\le L_p,\qquad
 v_p\binom{2h}h\le L_p^0,\qquad
 v_p\big(2r\binom{2r}r\big)\le\lfloor\log_p(2r)\rfloor.
\end{equation}
For the first identity write $a=h+1$, $b=n-h-1$ and use
$\binom{h+1/2}{n}=(-1)^b\binom{2a}a\binom{2b}b/(4^n\binom na)$.
At each odd prime-power level the two doubling carries minus the addition carry are at most1; the odd factors in the defining falling product have absolute value at most $2n-1$. For the last identity, a prime-power dividing $r$ contributes no doubling carry, so the valuation contribution per level is at most1. The central-binomial bound is the ordinary doubling-carry count. These arguments also prove the common-denominator, not merely individual-entry, assertions.

Generalized half-integer binomial coefficients are dyadic. So are $D_n(j)/n!=2^n\binom{j/2}{n}$ and the coefficients of \eqref{inverse:dyadic}. Moreover
$n!/D_n'(2i)=(-1)^{N-i}n\binom Ni/2^N$ and
$1/d_{2j+1}(0)=-(2j+1)\binom{2j}j/(2\,16^j)$.
Finally $\Lambda_j(2i)=\binom{i-1/2}{j}\binom{n-i-1/2}{N-j}$ is dyadic. Thus the odd part of $F$ for unit $\beta$ input is dyadic, while for unit $\sigma$ input its coefficients have valuation at least $-L_p$.

For odd $F$, \eqref{inverse:Laurent} and \eqref{inverse:counts} give a loss at most $L_p$ for every odd $p$, with no exceptional prime $3$. Taylor expansion of a Laurent monomial at $1$ has integral coefficients. This improves to $L_p^0$: only when $2n-1=p^e$ is there a difference. In that case the only possible top-level pole in \eqref{inverse:Laurent} is $h=j=n-1$. But $F_{2n-1}$ is divisible by $2n-1$ in $\Z[1/2]$, by \eqref{inverse:particular}. That term is $p$-integral, and every other term loses at most $e-1$ powers. This proves part 1 of Theorem~\ref{inverse:main}.

For even monomials, the algebraic coefficients of $H_{n,m}$ have odd denominator dividing $\Omega_{2n-2}$: polynomial primitives use odd denominators at most $2n-3$, and the rational terms of \eqref{inverse:primitiveR} are controlled by \eqref{inverse:counts}. The angle coefficient is dyadic. To see precisely what happens to its local jets, write $x=1+y$ and
\[
 w(y)=\sqrt{\frac{1+y}{1-y/3}},\qquad
 \sqrt3\,\frac{d}{dy}\left(\arctan\sqrt{\frac{1+y}{3-y}}-\frac\pi6\right)
   =\frac1{2\sqrt{1+y}\sqrt{1-y/3}}.
\]
All coefficients on the right are integral away from $2$ and $3$. Integrating through Taylor degree $2n-1$ costs at most $L_p$. Also
$\sqrt3u^{2r+1}=3^{-r}w^{2r+1}$ and
$\sqrt3u/(1+u^2)^r=w(3-y)^r/4^r$.
Thus every retained jet of \eqref{inverse:evenprimitive} has valuation at least $-L_p$ for $p\ge5$. Multiplication by $t^{1-2n}\sqrt{3/(4-t^2)}$, substitution $x=t^2$, Taylor truncation, and \eqref{inverse:dyadic} introduce no further such prime. The initial $F$ costs at most $L_p$, proving part 3.

For the pure $\tau$ construction only $1/\binom{2h}h$ introduces primes $p\ge5$; all other displayed normalized coefficients are dyadic or have powers of $3$ as denominators. This proves part 2.

For clarity we give the final $\alpha$ accounting rather than assuming cancellation. The recurrence for $d_j(1)$, $j\le2n-2$, gives $v_p(d_j(1))\ge-L_p$: the explicit coefficients are interlaced products
$-2\,4^r(j-1)(j-3)\cdots(j-2r+1)/(j(j-2)\cdots(j-2r))$, with at most one excess denominator multiple per odd prime-power level. In an off-diagonal entry of $q_0=t^{2i}/c_{2i}$, its $\tau$ coordinate therefore has valuation at least $-3L_p$, and its $\sigma$ coordinate at least $-2L_p$.
On the diagonal the normalized recurrence is
\[
 I_{2i}/c_{2i}=V+\sum_{\substack{2\le s\le2i\\s\ {
m even}}}
                          \frac{-4J_{s-2}+J_s}{s c_s}.
\]
Its $\tau$ coordinate is again at least $-3L_p$. Its $\sigma$ coordinate is
$-(H_{2i}-H_i)/3$, since $4c_{s-2}/c_s=s/(s-1)$, and loses at most $L_p$.
The correction by $\tau$ inverse columns costs at most one further $L_p$, and that by $\sigma$ columns at most two. Both totals are at most $4L_p$; $q_0$ itself loses at most $L_p$. This proves part 4 and the theorem.



% New improvement: closes the factor-three target-column denominator question left in Meiji R62.
\section{Three short-lcm powers suffice for the target inverse columns}
\label{inv3:sec:inverse-three}

Let $\mathcal M_n$ be the complete formal mixed-moment matrix with even
row exponents $0,2,\ldots,2n-2$, column exponents $0,\ldots,4n-1$, and
coordinate blocks $(\beta,\tau,\sigma,\alpha)$ for
$(1,\sqrt3,\pi,V)$. The Laurent and centered-angle construction gives
its inverse. For $p\ge5$, put
\[
 L=\lfloor\log_p(2n-1)\rfloor,\qquad C_a=\binom{2a}{a}.
\]
The already established inverse-column bounds are $-L$ for a
$\tau$ input and $-2L$ for a $\sigma$ input. The source Meiji R62~\cite{M62}
proved $-4L$ for the target input and explicitly left the apparent
improvement to $-3L$ unproved. We now establish that improvement.

\begin{theorem}[Improved target-column denominator]\label{inv3:thm:inverse-three}
Every entry of a unit $\alpha$ column of $\mathcal M_n^{-1}$ has
$p$-adic valuation at least $-3\lfloor\log_p(2n-1)\rfloor$
for every prime $p\ge5$. Consequently, the common inverse denominator
away from $2$ and $3$ divides
\[
 \operatorname{lcm}(1,\ldots,2n-1)^3.
\]
The sharper earlier bounds for the other three input blocks remain
valid. No optimality of the exponent three is asserted.
\end{theorem}

\begin{lemma}[Shared carries]\label{inv3:lem:shared-carries}
For distinct $0\le a,r<n$, the following valuations are at least $-L$:
\[
 \frac{C_r}{(a-r)C_a},\qquad
 \frac{C_r-C_a}{(a-r)C_a}.
\]
Also $v_p(sC_s)\le L$ for $1\le s<n$.
\end{lemma}
\begin{proof}
For every $h\le n-1$, Legendre's factorial formula gives
\[
 v_p(C_h)=\sum_{j=1}^{L}\epsilon_j(h),\qquad
 \epsilon_j(h)=\lfloor2h/p^j\rfloor-2\lfloor h/p^j\rfloor
              \in\{0,1\}.
\]
If $k=v_p(a-r)$, then $k\le L$ and
$\epsilon_j(a)=\epsilon_j(r)$ for $j\le k$.
Write $B$ for the sum of these common terms. Thus
$v_p(C_a)\le B+L-k$ and $v_p(C_r)\ge B$; also
$v_p(C_r-C_a)\ge B$. Division by $(a-r)C_a$ proves both bounds.
If $k=v_p(s)$, the first $k$ doubling carries of $s$ are zero.
Hence $v_p(C_s)\le L-k$, proving the last assertion.
\end{proof}

Write $d_{2a}(x)$ for the polynomial coefficient of the square-root
period in the usual conic moment recurrence:
\[
 d_0(x)=0,\qquad
 d_{2a}(x)=\frac{4(2a-1)}{2a}d_{2a-2}(x)-\frac{x^{2a-1}}a.
\]
Induction, using $C_a/C_{a-1}=2(2a-1)/a$, gives the exact identity
\begin{equation}\label{inv3:eq:d-sum}
 d_{2a}(x)=-C_a\sum_{s=1}^a\frac{x^{2s-1}}{sC_s}.
\end{equation}
In particular $v_p(d_{2a}(1))\ge v_p(C_a)-L\ge-L$.

\begin{lemma}[Normalized nuisance bounds]\label{inv3:lem:normalized-nuisance}
Let $q_0=t^{2a}/C_a$, where $0\le a<n$. For every even row
$2r$, $0\le r<n$, its square-root and pi coordinates satisfy
\[
 v_p(\tau_r(q_0))\ge-2L,\qquad
 v_p(\sigma_r(q_0))\ge-L.
\]
Its rational coordinates vanish, and its target coordinates are the
unit vector in position $a$.
\end{lemma}
\begin{proof}
For $r\ne a$, apart from $p$-adic units $2$ and $3$, the pi
coordinate is
$(C_a-C_r)/((r-a)C_a)$, so Lemma~\ref{inv3:lem:shared-carries} applies.
The square-root coordinate is
\[
 \frac{d_{2a}(1)-d_{2r}(1)}{2(r-a)C_a}.
\]
In the polynomial version of this quotient, each coefficient with
$s\le\min(a,r)$ equals, up to sign and the unit $2$,
\[
 \frac{C_r-C_a}{(a-r)C_a}\,\frac1{sC_s}.
\]
Each coefficient with $a<s\le r$ equals
$C_r/((a-r)C_a sC_s)$, up to a unit; each coefficient with
$r<s\le a$ equals $1/((a-r)sC_s)$, up to a unit.
The carry lemma bounds the first factor in every case below by $-L$,
and $v_p(sC_s)\le L$ bounds the second below by $-L$.
All polynomial coefficients therefore have valuation at least $-2L$,
as does their sum at $x=1$.

On the diagonal, the normalized recurrence reads
\[
 \frac{I_{2a}}{C_a}
 =V+\sum_{s=1}^a\frac{-4J_{2s-2}+J_{2s}}{2sC_s}.
\]
Its pi coefficient is $-(H_{2a}-H_a)/3$, of valuation at least $-L$.
Its square-root coefficient is
\begin{equation}\label{inv3:eq:diagonal-tau}
 -\sum_{s=1}^a\frac{2d_{2s-2}(1)+1}{2s^2C_s}.
\end{equation}
If $p\nmid s$, each summand has valuation at least $-2L$, by
\eqref{inv3:eq:d-sum} and $v_p(C_s)\le L$.
If $k=v_p(s)>0$, put $c=v_p(C_s)$. The identity
$C_s/C_{s-1}=2(2s-1)/s$ gives
$v_p(C_{s-1})=c+k$, while the carry lemma gives $c+k\le L$.
Equation \eqref{inv3:eq:d-sum} implies
$v_p(d_{2s-2}(1))\ge c+k-L$. This last number is nonpositive,
so the numerator of the summand in \eqref{inv3:eq:diagonal-tau} has
valuation at least $c+k-L$. After division, its valuation is
at least
\[
 (c+k-L)-(2k+c)=-L-k\ge-2L.
\]
This proves the diagonal bound. The other formal coordinates follow
directly from the mixed-moment recurrence: even $J$ moments have no
rational or target coordinate, and only the matching diagonal has
 target coefficient $C_a$.
\end{proof}

\begin{proof}[Proof of Theorem~\ref{inv3:thm:inverse-three}]
Start the unit target inverse column with $q_0=t^{2a}/C_a$.
Its coefficient valuation is at least $-L$. Subtract its nuisance
outputs using the established pure $\tau$ and pure $\sigma$ inverse
columns. Lemma~\ref{inv3:lem:normalized-nuisance} costs at most $2L$
and $L$ for these outputs; the respective inverse columns cost at most
$L$ and $2L$. Both corrections therefore have all coefficient
valuations at least $-3L$. Their finite sum has the same bound.
There is no rational-output correction. This proves the target-column
assertion and, together with the earlier block bounds, the common
short-lcm denominator statement.
\end{proof}

\paragraph{Scope.}
This closes a specific open denominator question in the supplied corpus.
It improves an exact rational inverse; clearing its denominator still
multiplies the target linear forms. It does not produce the independent
small integral controls required for an irrationality proof.



\section{Sharp local period images and denominator thresholds}
\subsection{Canonical coordinates and statements}
Let
\[
 \B(q,P)=\int_0^1\!\int_0^1\frac{q(s)P(t)}{\sqrt{1-s^2t^2/4}}\,ds\,dt,
 \qquad V=\B(1,1)=\operatorname{Cl}_2(\pi/3).
\]
For even $q$, write canonically
\[
 \B(q,P)=\beta+\tau\sqrt3+\sigma\pi+\alpha V.
\]
This means the rational coordinates defined by the recurrences below, not an assumption of linear independence of these constants. Write $\Rp$ for the rationals integral at $p$.

\begin{theorem}[Universal degree thresholds]\label{local:universal}
Let $p\ge5$ be prime, $q\in\Rp[s]$ even with $\deg q<p$, and $P\in\Rp[t]$.
\begin{enumerate}
\item If $\deg P<2p$ and $\tau=0$, then $\beta\in\Rp$.
\item If $\deg P<3p$ and $\tau=\sigma=0$, then $p\beta\in\Rp$.
\end{enumerate}
In both cases $\alpha\in\Rp$. Neither assertion requires orthogonality, positivity, a Rodrigues representation, or any restriction on the coefficients beyond local integrality.
\end{theorem}
The first assertion was proved in round 28; it is restated and proved here. The second is the new degree specialization of round 13's moment law.

\begin{theorem}[Exact unrestricted local image]\label{local:phase}
Put $q_m(s)=(1-s^2)^m$, where $m\ge0$ and $2m<p$. For $3\le M<p^2$, define
\[
 \Gamma_{m,M,p}=\{(\beta(P),\alpha(P)):P\in\Rp[t]_{\le M},\ \tau(P)=\sigma(P)=0\}.
\]
Then
\[
 \Gamma_{m,M,p}=\begin{cases}
 \Rp^2,&M<2p,\\
 p^{-1}\Rp\times\Rp,&2p\le M<3p,\\
 p^{-2}\Rp\times\Rp,&3p\le M<p^2.
 \end{cases}
\]
In particular both thresholds of Theorem~\ref{local:universal} are sharp, even for $m=0$.
\end{theorem}
This is an unrestricted polynomial image theorem. It does not assert that an orthogonal or positive sublattice attains these images.

\subsection{Moment formulas and the two pole layers}
Set $f(t)=\sqrt{1-t^2/4}$ and
\[
 J_j=\int_0^1t^j/f(t)\,dt,\qquad I_j=\int_0^1-t^j\log t/f(t)\,dt.
\]
In the coordinate order $(1,\sqrt3,\pi,V,\Lambda)$, where $\Lambda=\log((2+\sqrt3)/4)$, the initial values and recurrences are
\[
 J_0=(0,0,1/3,0,0),\quad J_1=(4,-2,0,0,0),\quad
 I_0=(0,0,0,1,0),\quad I_1=(4,-2,0,0,4),
\]
\[
 J_j=\frac{4(j-1)}jJ_{j-2}-\frac2j e_{\sqrt3},\qquad
 I_j=\frac{4(j-1)}jI_{j-2}-\frac4jJ_{j-2}+\frac1jJ_j.
\]
The monomial moment is
\[
 M_{ij}=\B(s^i,t^j)=\begin{cases}(J_j-J_i)/(i-j),&i\ne j,\\I_i,&i=j.\end{cases}
\]
All these identities follow by integration by parts. Even first indices have zero logarithmic coordinate. The $V$ coordinate is zero off the diagonal and equals $c_i=\binom{i}{i/2}$ on an even diagonal, proving the integrality assertion for $\alpha$.

We recall, with proof, the needed special case of round 13. Put $\epsilon=(-1/p)$ and $\eta=(-3/p)$, where the parentheses denote Legendre symbols.
\begin{lemma}[Low-first-index descent]\label{local:descent}
For even $0\le i<p$ and $j=ap+r<p^2$, $0\le r<p$,
\[
 (p^2M_{ij}^{\beta},p^2M_{ij}^{\tau},pM_{ij}^{\sigma})
 \equiv\begin{cases}c_i(\epsilon M_{0a}^{\beta},\eta M_{0a}^{\tau},M_{0a}^{\sigma}),&r=i,\\(0,0,0),&r\ne i\end{cases}\pmod p.
\]
All scaled coordinates displayed here are $p$-integral.
\end{lemma}
\begin{proof}
Use formal series
\[
 A=(1-4z^2)^{-1/2}=\sum c_jz^j,\quad U=\int_0^zA(u)du,\quad
 T=\int_0^z\frac{A(u)}{1-u}du,
\]
\[
 H=\tfrac13\int_0^z\frac{1-A(u)}u du,\quad U_2=\int_0^zU(u)du/u,\quad T_2=\int_0^zT(u)du/u.
\]
The recurrences give the generating functions
\[
 \sum J_j^\beta z^j=4AU,\quad \sum J_j^\tau z^j=-2AT,\quad \sum J_j^\sigma z^j=A/3,
\]
\[
 \sum I_j^\beta z^j=4AU_2,\quad\sum I_j^\tau z^j=-2AT_2,\quad\sum I_j^\sigma z^j=AH.
\]
In characteristic $p$, $A(z)=A_p(z)A(z^p)$ with $A_p=(1-4z^2)^{(p-1)/2}$, last coefficient $\epsilon$ and $A_p(1)=\eta$. The coefficients of $U,T,U_2,T_2,H$ at $j>0$ are respectively
\[
 c_{j-1}/j,\quad \Big(\sum_{k<j}c_k\Big)/j,\quad c_{j-1}/j^2,\quad
 \Big(\sum_{k<j}c_k\Big)/j^2,\quad-c_j/(3j).
\]
Through degree $<p^2$, only $j=ap$ survives after scaling by $p$ or $p^2$ respectively. The factorization of $A$ therefore gives
\[
 pU\equiv\epsilon U(z^p),\quad pT\equiv\eta T(z^p),\quad
 p^2U_2\equiv\epsilon U_2(z^p),\quad p^2T_2\equiv\eta T_2(z^p),\quad pH\equiv H(z^p).
\]
Multiplication by $A$ proves the corresponding $J,I$ residues and valuation bounds. If $r\ne i$, the moment divisor $i-j$ is a unit and the extra $p$ kills the residue. If $r=i$ and $a>0$, divide the $J$ difference by $-ap$. The case $a=0,r=i$ is a low integral diagonal and both sides are zero. This proves the lemma.
\end{proof}

The three high columns, with rows $(\beta,\tau,\sigma)$, are
\begin{equation}\label{local:small}
 W_1=M_{01}=(-4,2,1/3)^T,\quad W_2=M_{02}=(0,1/2,-1/6)^T,\quad
 W_3=M_{03}=(-32/9,2,1/9)^T.
\end{equation}
Their first two nuisance columns have determinant $-1/2$, while
\begin{equation}\label{local:dets}
 \det[W_1,W_2,W_3]=2/9,\qquad
 \beta=-\tfrac43\tau-4\sigma\quad\hbox{on }\operatorname{span}(W_1,W_2).
\end{equation}
For $\deg P<3p$, set $\phi_a=\sum_{i\text{ even}<p}q_iP_{ap+i}c_i$, $a=0,1,2$. Lemma~\ref{local:descent} says that the scaled coordinates descend to $\sum\phi_aM_{0a}$. The column $M_{00}$ has only a $V$ coordinate. Thus $\tau=\sigma=0$ forces $\phi_1=\phi_2=0$, and then $p^2\beta\equiv0$. This proves the second assertion of Theorem~\ref{local:universal}. More explicitly, throughout this degree range,
\begin{equation}\label{local:congruence}
 p\beta+\frac{4\epsilon}{3\eta}\,p\tau+4\epsilon\sigma\in\Rp.
\end{equation}

\subsection{The first threshold and low-degree surjectivity}
For completeness we give the stronger first-threshold argument from round 28. Put $c_p=J_p^\beta/J_p^\tau$. The odd antiderivative at $p=2h+1$ gives
\[
 J_p=2\sum_{r=0}^{h}(-1)^r\binom hr4^{h-r}\frac{2^{2r+1}-3^r\sqrt3}{2r+1}.
\]
Only $r=h$ has a $p$ denominator, so $pJ_p^\beta\equiv4\epsilon$ and $pJ_p^\tau\equiv-2\eta\ne0$. In particular $c_p$ is a unit. The recurrence for
$A_j=J_j^\beta-c_pJ_j^\tau$ is integral for $j<2p$, with its only singular step removed by $A_p=A_0=0$. Comparing recurrence steps $p+i$ and $i$ gives
\[
                         A_{p+i}\equiv A_i\pmod p\qquad(i<p\text{ even}).
\]
The only off-diagonal moment divisor divisible by $p$ is $j=p+i$, so $\beta-c_p\tau$ is integral on every even $i<p,j<2p$. Low diagonals are integral. This proves the first assertion.

We also need a low-degree surjectivity fact, again restating round 28's explicit boundary model. Let $H_m(t)=\int_0^1q_m(s)/f(st)\,ds$, and put $h_m=H_m(1)$, $h'_m=H'_m(1)$. Its differential operator and adjoint are
\[
 L_m=t(4-t^2)D^2+(8m+8-3t^2)D-t,\quad
 L_m^*=t(4-t^2)D^2-(8m+3t^2)D-t.
\]
Integration by parts gives canonically
\begin{equation}\label{local:boundary}
 \B(q_m,L_m^*R)=[3R'(1)-(8m+4)R(1)]h_m-3R(1)h'_m+\kappa_mR(0),\quad
 \kappa_m=4^{m+1}/\binom{2m}m.
\end{equation}
The $V$ coefficient is checked directly by the binomial recurrence; the other three coordinates follow either from the moment recurrences or their known rational linear independence. The latter uses transcendence of $\pi$, not any unproved fact about $V$.
The initial values $h_0=\pi/3$, $h_1=\sqrt3-\pi/3$, $h'_0=2\sqrt3/3-\pi/3$ and
\[
 h'_m=2mh_{m-1}-(2m+1)h_m,\quad
 (m+1)h_{m+1}+(2m+1)h_m-3mh_{m-1}=0\quad(m>0)
\]
show that
\[
 \det\begin{pmatrix}h_m^\tau&h_m^\sigma\\(h'_m)^\tau&(h'_m)^\sigma\end{pmatrix}
                 =-\tfrac29(-3)^m.
\]
This and $\kappa_m$ are units when $2m<p$. Taking $R=(1-t)^2$ in \eqref{local:boundary} represents a rational unit. Quadratic Hermite interpolation with $R(0)=0$ independently supplies arbitrary integral nuisance coordinates. Finally $\B(q_m,1)$ has $\alpha=1,\beta=0$; its nuisance coordinates are integral and can be cancelled in this way. Consequently the degree-three output map, with all four coordinates, is an isomorphism over $\Rp$. Its determinant in row order $(\beta,\alpha,\tau,\sigma)$ is explicitly
\[
                        -\frac{2\,4^m(-3)^m}{9\binom{2m}m},
\]
a unit. Thus degree-three polynomials realize every integral formal output.

\subsection{Sharpness and proof of the local phase theorem}
Write $C_j$ for the full output of $\B(q_m,t^j)$, and retain only its three coordinates $(\beta,\tau,\sigma)$ when applying determinants. Since $q_m(0)=1$ and $\deg q_m<p$, Lemma~\ref{local:descent} gives
\begin{equation}\label{local:high}
 \operatorname{diag}(p^2,p^2,p)[C_p,C_{2p},C_{3p}]
       \equiv\operatorname{diag}(\epsilon,\eta,1)[W_1,W_2,W_3]\pmod p.
\end{equation}
All three columns have $\alpha=0$.
For degree $2p$, the quotient $x=C_{2p}^\tau/C_p^\tau$ is integral and satisfies $x\equiv1/4\pmod p$. The polynomial
\[
                           P_*=p(t^{2p}-x t^p)
\]
has $\tau=0$, integral $\sigma$, and $\alpha=0$. Moreover
\[
                           p\beta(P_*)\equiv\epsilon\ne0\pmod p.
\]
Degree-three surjectivity supplies a correction with integral coefficients and output $(0,0,0,\sigma(P_*))$ in row order $(\beta,\alpha,\tau,\sigma)$. Subtract it. This gives a nuisance-free simple pole with zero $\alpha$, proving attainment in the middle range.

For degree $3p$, solve the two scaled nuisance equations using $C_p,C_{2p}$ to cancel those of $C_{3p}$. Their scaled determinant is $-\eta/2$, a unit. The resulting coefficients are in $\Rp$, and the full determinant in \eqref{local:high} is $2\epsilon\eta/9$, also a unit. Hence the resulting nuisance-free output has $p^2\beta$ a unit and $\alpha=0$.
Together with the integral pair generators, these constructions give every claimed pair. The reverse inclusions follow from Theorem~\ref{local:universal} in the first two ranges, and the pole-order bounds in Lemma~\ref{local:descent} in the third. This proves Theorem~\ref{local:phase}.
All witnesses can be made genuinely integral polynomials by multiplying by a common denominator prime to $p$; their exact pole orders are unchanged.

\subsection{A common denominator and the improved Bernstein budget}
Write $D_N=\operatorname{lcm}(1,\ldots,N)$. For integral input polynomials of degrees at most $N$, the rational and square-root moment coordinates have denominator dividing $4D_N^2$. Here is a short verification. For odd and even indices respectively, the incomplete primitive $J_j(t)=\int_0^tu^j/f(u)du$ is
\[
 J_{2r+1}(t)=2^{2r+2}\sum_{k=0}^r\frac{(-1)^k\binom rk}{2k+1}(1-f(t)^{2k+1}),
\]
\[
 J_{2r}(t)=\binom{2r}rJ_0(t)+2tf(t)\sum_{k=1}^r\frac{(-1)^k\binom{2r}{r-k}}kU_{k-1}(1-t^2/2).
\]
The polynomial $U_{k-1}(1-z/2)$ is integral. Expanding these formulas gives coefficients with denominator dividing $D_j$. Now $I_j=\int_0^1J_j(t)dt/t$, with
$\int_0^1(f(t)-1)dt/t=-\Lambda-1+\sqrt3/2$ and
$\int_0^1t^a f(t)dt=J_a-J_{a+2}/4$; all resulting indices are at most $j$. This pays at most one further $D_j$ and the factor $4$. The off-diagonal moment divisor is already paid by the second $D_N$.

\begin{corollary}[Explicit general denominator]\label{local:denominator}
Let $q,P$ be integral, $q$ even, $\deg q\le D$, $\deg P\le M$, and $\tau=\sigma=0$. Put $N=\max(D,M,1)$ and
\[
 u=\max(3,D,\lfloor M/2\rfloor),\qquad v=\max(3,D,\lfloor M/3\rfloor),
\]
\[
 G_{D,M}=\frac{4D_N^2}{\displaystyle\prod_{u<p\le N}p^2\prod_{v<p\le\min(u,N)}p}.
\]
Products run over primes. Then $G_{D,M}\beta\in\Z$ and $\alpha\in\Z$. If $D/n\to\delta$, $M/n\to\mu$ with positive finite $\max(\delta,\mu)$, then
\[
 \frac{\log G_{D,M}}n\longrightarrow
                \max(\delta,\mu/2)+\max(\delta,\mu/3).
\]
\end{corollary}
\begin{proof}
Every removed prime exceeds $D$ and $3$. In the first product $M<2p$, so all of its denominator is removed. In the second $M<3p$, so at most a simple pole remains. Also $N<3p<p^2$ for every removed prime, so its original exponent in $D_N$ is one and the displayed quotient is an integer. All other prime powers retain the global bound. The prime number theorem gives the asserted limit.
\end{proof}



% New synthesis of Taisho R36--R38, not a claim of irrationality.
% Requires amsmath, amssymb, amsthm, theorem/corollary/proof environments.
\section{Optimizing the Bernstein degree and denominator budget}
\label{opt:taisho:optimization}

The positive-series and saddle formulas are from Taisho R36. Combining
them with the stronger universal degree denominator of Taisho R38 admits
an optimization which was not performed in those source rounds. It distinguishes the best analytic and
arithmetic balance from the unsolved problem of realizing that balance by
positive integral controls. In particular, the uniform positive cone proved
at $m/n=3/20$ is not assumed to persist at the new optimal parameter.

Let $n,m,K$ be nonnegative integers with $n>0$, put
\[
 q_m(s)=(1-s^2)^m,\qquad
 P_{n,K,d}(t)=\frac{(-1)^n}{n!}\frac{d^n}{dt^n}
          \left[t^{n+d}(1-t)^{n+K-d}\right],\quad 0\le d\le K,
\]
and define $T_{n,m,K,d}=\mathcal B(q_m,P_{n,K,d})$.
Every $P_{n,K,d}$ has integer coefficients. Write
$T_{n,m,K}=T_{n,m,K,K}$ for the last column.
For fixed real parameters $\rho>0$ and $\kappa\ge0$, define
\begin{align}
 F_{\rho,\kappa}(x)={}&-x\log2+
 \frac{x}{2}\log\frac{x}{x+2\rho}
 +\rho\log\frac{2\rho}{x+2\rho}
 +x\log x-(x-1)\log(x-1)\notag\\
 &+(x+\kappa)\log(x+\kappa)
 -(x+\kappa+1)\log(x+\kappa+1),\qquad x>1,\label{opt:taisho:F}\\
 \chi(\rho,\kappa)&=-\max_{x>1}F_{\rho,\kappa}(x),\label{opt:taisho:chi}\\
 g(\rho,\kappa)&=
 \max\left(2\rho,\frac{1+\kappa}{2}\right)
 +\max\left(2\rho,\frac{1+\kappa}{3}\right),\label{opt:taisho:g}\\
 B(\rho,\kappa)&=\chi(\rho,\kappa)-g(\rho,\kappa).
 \label{opt:taisho:budget}
\end{align}

\begin{theorem}[A sharp budget optimization within this family]
\label{opt:taisho:main}
The following statements hold.
\begin{enumerate}
\item If $m/n\to\rho>0$ and $K/n\to\kappa\ge0$, then
\[
 0<T_{n,m,K,d}\le T_{n,m,K}\quad(0\le d\le K),\qquad
 \log T_{n,m,K}=-\chi(\rho,\kappa)n+o(n).
\]
When $m=\rho n+O(1)$ and $K=\kappa n+O(1)$, the last error is
$O_{\rho,\kappa}(\log n)$.
\item For each fixed $\kappa\ge0$, the function
$\rho\mapsto B(\rho,\kappa)$ has a unique maximum, at
\[
                       \rho=\frac{1+\kappa}{6}.
\]
\item The optimized function
\[
 B_*(\kappa)=B\left(\frac{1+\kappa}{6},\kappa\right)
\]
is concave and strictly decreasing on $[0,\infty)$. Its endpoint value is
\[
 B_*(0)=1.351404651210758218\ldots.
\]
In particular, for full width $K=n$, the optimal binomial order is
$m=n/3+O(1)$ and
\[
 B_*(1)=1.175542931127682114\ldots.
\]
These maxima concern the explicitly defined budget
\eqref{opt:taisho:budget}, not all possible denominator estimates or
approximation constructions.
\end{enumerate}
\end{theorem}

\begin{proof}
Integrating Rodrigues' formula by parts, and using the positive kernel
series, gives the exact expansion
\begin{equation}\label{opt:taisho:positive-series}
 T_{n,m,K,d}=
 \sum_{r\ge\lceil n/2\rceil}
 \frac{\binom{2r}{r}}{16^r}
 \frac{4^m(2r)!m!(r+m)!}{r!(2r+2m+1)!}
 \frac{(2r)!(2r+d)!(n+K-d)!}
      {n!(2r-n)!(2r+n+K+1)!}.
\end{equation}
All terms are strictly positive. At fixed $r$, the $d$-dependent factor
$(2r+d)!(n+K-d)!$ is log-convex. Its value at $d=K$ is at least its value
at $d=0$, because
\[
 \frac{(2r+K)!n!}{(2r)!(n+K)!}
       =\prod_{j=1}^K\frac{2r+j}{n+j}\ge1.
\]
The largest endpoint therefore bounds every column term by term.

With $x=2r/n$, Stirling's bounds applied to the last-column summand in
\eqref{opt:taisho:positive-series} give $nF_{\rho,\kappa}(x)+O(\log n)$
on each compact subinterval of $x>1$, when the parameter errors are
$O(1)$. The derivatives needed to control its maximum are
\begin{align*}
 F_x&=-\log2+\frac12\log\frac{x}{x+2\rho}
             +\log\frac{x}{x-1}+\log\frac{x+\kappa}{x+\kappa+1},\\
 F_{xx}&=\frac{\rho}{x(x+2\rho)}
       -\frac1{x(x-1)}+
        \frac1{(x+\kappa)(x+\kappa+1)}.
\end{align*}
For $1<x<2$, the second expression is bounded above by
\[
 \frac1{2x}-\frac1{x(x-1)}+\frac1{x(x+1)}
              =\frac{x^2-5}{2x(x^2-1)}<0.
\]
Also $F_x\to+\infty$ as $x\downarrow1$. For $x\ge2$,
\[
 F_x< -\log2+\log\frac{x}{x-1}\le0.
\]
Thus there is exactly one maximizer $x_{\rho,\kappa}\in(1,2)$,
it is nondegenerate, and it is determined by
\begin{equation}\label{opt:taisho:saddle}
 x^3(x+\kappa)^2
   =4(x+2\rho)(x-1)^2(x+\kappa+1)^2.
\end{equation}

For completeness, the tail of the positive sum needs no unproved uniform
saddle assertion. The ratio of consecutive last-column summands is a
rational function of $r,n,m,K$. As $n\to\infty$ with $x=2r/n$ fixed,
it converges to $\exp(2F_x(x))$. On $x\ge2$ this is bounded by a constant
strictly less than one for the fixed parameters; the bound also holds at
infinity, where the ratio tends to $1/4$. Consequently the same geometric
bound holds for the exact ratios for all sufficiently large $n$ and
$r\ge n$. On the remaining compact interval, the unique nondegenerate
maximum and Stirling's inequalities give the usual two-sided positive
Laplace bounds. This proves the asserted logarithmic asymptotic. Mere
convergence of the degree ratios gives an $o(n)$ error by continuity,
with the same bounds on compact parameter neighborhoods.

The denominator theorem for degrees $D=2m$ and $M=n+K$ gives
\[
 \log G_{2m,n+K}=g(\rho,\kappa)n+o(n).
\]
There is no claim here that this universal denominator is the exact
primitive denominator on the orthogonal or positive sublattice.
Smooth dependence of the unique saddle and the envelope rule give
\begin{equation}\label{opt:taisho:envelope}
 \chi_\rho=\log\left(1+\frac{x_{\rho,\kappa}}{2\rho}\right)>0,
 \qquad
 \chi_\kappa=\log\left(1+\frac1{x_{\rho,\kappa}+\kappa}\right).
\end{equation}
Below $\rho=(1+\kappa)/6$, the function $g$ is the constant
$5(1+\kappa)/6$, so the budget strictly increases.
At and above this corner, $g$ has slope at least $2$ with respect to
$\rho$, while
\[
 \chi_\rho<\log\left(1+\frac{6}{1+\kappa}\right)
          \le\log7<2.
\]
The budget strictly decreases above the corner. This proves its unique
maximization in $\rho$.

For fixed $x$, the Hessian of $F$ in $(\rho,\kappa)$ is diagonal with
strictly positive entries
\[
 F_{\rho\rho}=\frac{x}{\rho(x+2\rho)},\qquad
 F_{\kappa\kappa}=\frac1{(x+\kappa)(x+\kappa+1)}.
\]
A pointwise supremum of convex functions is convex. It follows that
$\chi=-\sup_xF$ is concave jointly in $(\rho,\kappa)$, and hence
\[
 B_*(\kappa)=\chi\left((1+\kappa)/6,\kappa\right)
                   -5(1+\kappa)/6
\]
is concave. By \eqref{opt:taisho:envelope}, its right derivative at zero is
\begin{equation}\label{opt:taisho:derivative}
 B_*'(0+)=\frac16\log(1+3x_0)
             +\log(1+1/x_0)-\frac56,
\end{equation}
where $x_0$ is the unique root in $(1,2)$ of
\[
                  3x^5=4(3x+1)(x-1)^2(x+1)^2.
\]
The exact rational enclosure
\[
 1.34548950909964750934<x_0<1.34548950909964750935
\]
and elementary rational logarithm bounds give
\[
               -0.008147<B_*'(0+)<-0.008145<0.
\]
Concavity now implies strict decrease on the entire half-line.
The two displayed budget constants follow from the same interval
calculation, using \eqref{opt:taisho:F}. At $\kappa=1$ the saddle is enclosed by
\[
 1.41091690957728730184<x_{1/3,1}
                          <1.41091690957728730185.
\]
All the inequalities can be replayed with rational arithmetic. For
example, reduce a positive rational to $y\in[1,2]$, set
$z=(y-1)/(y+1)$, and bound the logarithm by
\[
 2\sum_{j=0}^{N-1}\frac{z^{2j+1}}{2j+1}
 \le\log y\le
 2\sum_{j=0}^{N-1}\frac{z^{2j+1}}{2j+1}
 +\frac{2z^{2N+1}}{(2N+1)(1-z^2)}.
\]
This certifies the signs and decimal enclosures independently of
floating-point optimization.
\end{proof}

\subsection{Including every shifted binomial first factor}
The older shifted family of Taisho R24 allows
$q_{a,m}(s)=s^{2a}(1-s^2)^m$. Let $a/n\to\eta\ge0$ and
$m/n\to\rho\ge0$. With $0\log0=0$, replace the first-factor part of
\eqref{opt:taisho:F} by
\[
 h(x/2+\eta,\rho),\qquad
 h(y,r)=y\log y+r\log r-(y+r)\log(y+r).
\]
More precisely, put
\begin{align*}
 \widetilde F_{\eta,\rho,\kappa}(x)
 ={}&-x\log2+h(x/2+\eta,\rho)
 +x\log x-(x-1)\log(x-1)\\
 &+(x+\kappa)\log(x+\kappa)
 -(x+\kappa+1)\log(x+\kappa+1),\\
 \widetilde\chi(\eta,\rho,\kappa)
 ={}&-\max_{x>1}\widetilde F_{\eta,\rho,\kappa}(x),\\
 \widetilde g(\eta,\rho,\kappa)
 ={}&\max\left(2(\eta+\rho),\frac{1+\kappa}{2}\right)
 +\max\left(2(\eta+\rho),\frac{1+\kappa}{3}\right),\\
 \widetilde B={}&\widetilde\chi-\widetilde g.
\end{align*}
The same positive-series argument applies: its first-factor moment is now
\[
 \int_0^1s^{2r+2a}(1-s^2)^m\,ds
 =\frac12\mathrm B(r+a+1/2,m+1).
\]
Thus $\widetilde B$ is precisely the analogue of the preceding
universal-denominator budget for the entire shifted family.

\begin{theorem}[Global optimum with shifts included]
\label{opt:taisho:shiftglobal}
For every $\eta,\rho,\kappa\ge0$,
\[
 \widetilde B(\eta,\rho,\kappa)
 \le B_*(0)-\log(3x_0)\eta+B_*'(0+)\kappa.
\]
Consequently the unique global maximum of this budget is
$B_*(0)=1.351404651210758218\ldots$, attained at
$(\eta,\rho,\kappa)=(0,1/6,0)$.
For fixed full width $\kappa=1$,
\[
 \widetilde B(\eta,\rho,1)
 \le B_*(1)+\log\left(\frac{2}{3x_1}\right)\eta,
 \qquad x_1=x_{1/3,1},
\]
and its unique maximum is attained at $(\eta,\rho)=(0,1/3)$.
\end{theorem}
\begin{proof}
The entropy function $h(y,r)$ is convex on $y>0,r\ge0$.
For $r>0$ this follows from
\[
 \nabla^2h=\frac1{y+r}
 \begin{pmatrix}r/y&-1\\-1&y/r\end{pmatrix},
\]
a positive semidefinite matrix, and the assertion at $r=0$ follows by
continuity. The remaining dependence of $\widetilde F$ on $\kappa$
is convex, as before. Thus $\widetilde\chi$ is jointly concave;
$\widetilde g$, a sum of maxima of affine functions, is convex.
Their difference is concave on the nonnegative parameter octant.

The unique nondegenerate saddle still lies in $(1,2)$: the first-factor
contribution to $F_x$ is now
$\tfrac12\log((x+2\eta)/(x+2\eta+2\rho))\le0$, and its contribution
to $F_{xx}$ is at most $1/(2x)$. The earlier derivative argument applies,
including $\rho=0$ by continuity. In particular the envelope derivatives
at either interior-in-$\rho$ point used below are legitimate.

At $(\eta,\rho,\kappa)=(0,1/6,0)$, set
\[
 A=\log(1+3x_0),\qquad C=\log(1+1/x_0).
\]
The gradient of $\widetilde\chi$ at that point is
\[
 \bigl(\log(1+1/(3x_0)),\ A,\ C\bigr).
\]
Since $0<A<\log7<2$, the vector
\[
                  (A,A,5/6-A/6)
\]
is a subgradient of $\widetilde g$ there. Indeed the first maximum in
$\widetilde g$ selects $(1+\kappa)/2$, and the tied second maximum
can be averaged with weight $A/2$ on $2(\eta+\rho)$.
The supporting-plane inequality for the concave difference therefore
gives
\[
 \widetilde B(\eta,\rho,\kappa)
 \le B_*(0)+\bigl[\log(1+1/(3x_0))-A\bigr]\eta
                 +\bigl[C-5/6+A/6\bigr]\kappa.
\]
The two bracketed coefficients are respectively $-\log(3x_0)$ and
$B_*'(0+)$. Both are strictly negative. If either $\eta>0$ or
$\kappa>0$, the value is strictly below $B_*(0)$. Along
$\eta=\kappa=0$, the unique optimum in $\rho$ was proved in
Theorem~\ref{opt:taisho:main}.

For $\kappa=1$, repeat the same supporting-plane calculation at
$(\eta,\rho)=(0,1/3)$, using
$A=\log(1+3x_1/2)\in(0,2)$. The $\rho$ coefficient is zero and the
$\eta$ coefficient is $\log(2/(3x_1))<0$. The earlier unique optimum
on $\eta=0$ completes the proof.
\end{proof}

The optimizer $\kappa=0$ is a boundary budget, not an existence theorem
for positive cancelling controls. In particular, the single-column
case $K=0$ does not thereby acquire the two missing nuisance
cancellations. Conversely, larger windows can improve the arithmetic
existence of controls even while worsening this particular budget.

\begin{corollary}[An improved full-width sufficient criterion]
\label{opt:taisho:newcriterion}
Suppose that, for infinitely many positive integers $n$ divisible by $3$,
there are rational numbers $w_0,\ldots,w_n\ge0$, not all zero, such that
\[
 P_w=\sum_{d=0}^n w_dP_{n,n,d}\in\mathbb Z[t],\qquad
 \tau(\mathcal B(q_{n/3},P_w))=
 \sigma(\mathcal B(q_{n/3},P_w))=0,
\]
and
\[
 \max_d w_d\le\exp(cn)
 \quad\hbox{for a fixed }c<1.175542931127682114\ldots.
\]
Then $V$ is irrational.
\end{corollary}
\begin{proof}
The universal denominator has rate $5n/3+o(n)$, and the target
coordinate $\alpha$ is integral because both input polynomials are
integral. The positive series and the endpoint estimate give
\[
 0<G_{2n/3,2n}\,\mathcal B(q_{n/3},P_w)
       \le\exp\bigl((c-B_*(1))n+o(n)\bigr)\longrightarrow0.
\]
If $V=a/b$ were rational in lowest terms, multiplying this positive
number by $b$ would give a positive integer tending to zero.
\end{proof}

The corollary permits rational Bernstein weights only when their actual
Rodrigues output is coefficientwise integral. It does not turn a bounded
positive rational kernel into an integral one without paying its true
arithmetic cost. The positive-cone theorem at $(\rho,\kappa)=(3/20,1)$
and its three-dimensional quotient bound therefore cannot simply be
transferred to $(1/3,1)$. The new conclusion closes the optimization of
this universal budget. It does not close the infinite control-height
problem or prove irrationality of $V$.


% Unified-paper section: new combination of Reiwa R49 and Keio R51.
% Needs amsmath/amssymb/amsthm and theorem/lemma/corollary environments.
\section{Complete first-layer saturation and the missing two-block interval}
\label{sat:reiwa:newsec}

Let $L_j(t)=P_j(2t-1)$ be the shifted Legendre polynomial, and put
\[
 \mathcal O_{n,K}=\{G\in\mathbb Q[t]:\deg G\le n+K,
       \ \int_0^1t^iG(t)\,dt=0\ (0\le i<n)\}.
\]
We distinguish the nested lattice and its coefficientwise saturation,
\[
 \Lambda=\sum_{j=0}^K\mathbb ZL_{n+j},\qquad
 \Lambda^{\rm sat}=\mathcal O_{n,K}\cap\mathbb Z[t].
\]
For a rational polynomial $P$, its canonical coordinates are specified by
\[
 P=P(0)+\frac t2\mathcal DU,\qquad
 \mathcal D=(4-t^2)\partial_t^2-3t\partial_t-1,
\]
\[
 C(P)=(\alpha,\beta,\tau,\sigma)
 =\left(P(0),2U(0),-U(1),\frac{3U'(1)-U(1)}6\right).
\]
The diagonal coefficients of $\mathcal D$ are $-(j+1)^2$, so this
specification is unambiguous. No independence of the associated real
periods is assumed. For a coefficient lattice $\mathcal L$, let
$\Gamma(\mathcal L)$ be the pair image $(\beta,\alpha)$ of $F\mathcal L$
subject to the exact equations $\tau=\sigma=0$.

The new point of this section is an \emph{upper} bound for the whole
saturated image. Selected divided columns give lower bounds, but do not
by themselves show that further integral columns cannot enlarge the
image. The complete reflected basis makes that distinction explicit.

\subsection{The complete coefficient-selected source space}
For a prime $p$, write
\[
 \mathcal C_pP(x)=\sum_{j\ge0}\overline{[t^{jp}]P(t)}x^j.
\]
The following lemma records the complete saturation mechanism from
Reiwa R49~\cite{R49}, with its proof for the present use.

\begin{lemma}\label{sat:reiwa:reflectionbasis}
Let $p>3$, $d=n+K$, and $2\lfloor d/p\rfloor+2<p$. For each reflected
pair $hp+l,hp+p-1-l\in[n,d]$, $0\le l<(p-1)/2$, replace the larger
nested column by
\[
 D_{h,l}=\frac{L_{hp+p-1-l}-L_{hp+l}}p.
\]
These columns form a basis of $\Lambda^{\rm sat}_{(p)}$. Moreover,
\[
 D_{h,l}\equiv(2h+1)L_h(t^p)D_{0,l}(t)\pmod p.                 \tag{\ref{sat:reiwa:reflectionbasis}.1}
\]
Here $D_{0,l}(0)=0$, $\deg D_{0,l}=p-1-l$, and its leading
coefficient is a $p$-unit.
\end{lemma}
\begin{proof}
The binomial expression
\[
 [t^j]L_N=(-1)^{N+j}\binom Nj\binom{N+j}j
\]
gives $L_{hp+r}\equiv L_h(t^p)L_r$ and
$L_{p-1-l}\equiv L_l$ modulo $p$. This proves integrality of the
divided columns. For completeness, put
$\mathcal L_l=t(1-t)\partial_t^2+(1-2t)\partial_t+l(l+1)$.
The Legendre equations give
\[
 \mathcal L_l D_{h,l}\equiv
 (2h+1)(2l+1)L_h(t^p)L_l,\qquad
 \mathcal L_lD_{0,l}\equiv(2l+1)L_l.
\]
Thus the difference in the asserted congruence is a polynomial
solution of $\mathcal L_lY=0$. Every such solution is
$L_l(t)A(t^p)$: the Wronskian equation gives
$(Y/L_l)'=B(t^p)/(t(1-t)L_l^2)$; if $B\ne0$, its leading Laurent
term at zero has exponent congruent to $-1$ modulo $p$, impossible
for a derivative. The roots of $L_l$ are simple, so no pole of a
rational function of $t^p$ can be cancelled by $L_l$.

The coefficients of degrees $jp$, $0\le j\le h$, fix $A$.
Grouping factorials in complete $p$-blocks gives
\[
 [t^{jp}]L_{hp+r}\equiv
 (-1)^{h+r+j}\binom hj\binom{h+j}j
       (1+2jpH_r)\pmod{p^2},\qquad H_r=\sum_{i=1}^r1/i.
\]
The block identity used here is
$\binom{ap}{bp}\equiv\binom ab\pmod{p^2}$ for $a<p$;
it follows by grouping nonmultiples of $p$ and using
$\sum_{i=1}^{p-1}1/i\equiv0\pmod p$.
Since $H_{p-1-l}\equiv H_l$ and the signs agree, all selected
coefficients of $D_{h,l}$ vanish. Those of
$L_h(t^p)D_{0,l}$ also vanish because $D_{0,l}(0)=0$ and its
degree is below $p$. As $L_l(0)$ is a unit, $A=0$.

An ordinary column has reduced leading degree
$hp+\min(r,p-1-r)$. Two such degrees coincide precisely for a
reflected pair. After division by $p$, the new column has reduced
leading degree $hp+p-1-l$: its leading binomial coefficient has
exactly one low-digit carry and no high carry. All basis columns
therefore have distinct reduced leading degrees and are independent
modulo $p$. Their lattice is $p$-saturated in its rational span:
reducing a relation $pG=\sum c_jB_j$ forces every $c_j$ to be a
multiple of $p$. This proves that the displayed lattice is the
entire coefficient saturation.
\end{proof}

\begin{theorem}[Complete first-layer source classification]
\label{sat:reiwa:classification}
Assume the hypotheses of Lemma~\ref{sat:reiwa:reflectionbasis}. Write
\[
 a=pu+s,\quad b=pv+w,\quad 0\le s,w<p,\quad m=s+w<p/2,
\]
\[
 F=(t-2)^a(t+2)^b,\qquad F_1(x)=(x-2)^u(x+2)^v.
\]
Let
\[
 \mathcal H=\{h:[hp,(h+1)p-1]\cap[n,d]\ne\varnothing\}
\]
and, if $m\ge1$, let
\[
 \mathcal J=\{h\in\mathcal H:hp+m-1\ge n,
                         \ (h+1)p-m\le d\}.
\]
Set $\mathcal J=\varnothing$ when $m=0$. Then the exact
$\mathbb F_p$-linear image is
\[
 \mathcal C_p(F\Lambda^{\rm sat}_{(p)})
 =\operatorname{span}_{\mathbb F_p}
   \left(\{F_1L_h:h\in\mathcal H\}
         \cup\{F_1xL_h:h\in\mathcal J\}\right).       \tag{\ref{sat:reiwa:classification}.1}
\]
The nested image is exactly the span of the first set alone.
\end{theorem}
\begin{proof}
Put $F_0=(t-2)^s(t+2)^w$. Frobenius gives
$F\equiv F_1(t^p)F_0(t)$. Every ordinary column in block $h$
reduces to $L_h(t^p)L_l(t)$ with $l\le(p-1)/2$.
The low product $F_0L_l$ has degree below $p$ and unit constant
coefficient. Its selected image is thus a nonzero scalar multiple
of $F_1L_h$. This proves the nested assertion.

For a divided basis column, Lemma~\ref{sat:reiwa:reflectionbasis}
gives the selected image of
$(2h+1)F_1(t^p)L_h(t^p)F_0D_{0,l}$.
The low product has zero constant coefficient and degree below
$2p$. Its selected polynomial is therefore $c_lx$ for a scalar
$c_l$. It is zero whenever $l\ge m$, because then its degree is
below $p$.

A reflected pair with $l<m$ is available only if
$hp+m-1\ge n$ and $(h+1)p-m\le d$. Conversely, whenever these
inequalities hold, the pair with $l=m-1$ is available. Its low
product has degree exactly $p$, with unit leading coefficient
\[
 \frac1p\binom{2(p-m)}{p-m};
\]
there is exactly one carry because $p/2<p-m<p$.
The factor $2h+1$ is a unit by the size hypothesis. Thus it
supplies $F_1xL_h$. The full basis lemma now proves both
inclusions in the claimed image, rather than only a lower bound.
\end{proof}

\begin{corollary}[Only boundary blocks enlarge the first-layer image]
\label{sat:reiwa:edges}
In Theorem~\ref{sat:reiwa:classification}, suppose
$\mathcal H=\{H,H+1,\ldots,J\}$ with $H<J$. Let
$\epsilon_-=1$ if $H\ge1$ and $H\in\mathcal J$, and zero
otherwise; let $\epsilon_+=1$ if $J\in\mathcal J$, and zero
otherwise. Then the saturated selected image is exactly
\[
 F_1\operatorname{span}_{\mathbb F_p}
       \{L_{H-\epsilon_-},\ldots,L_{J+\epsilon_+}\}.
\]
Its dimension exceeds the nested dimension by precisely
$\epsilon_-+\epsilon_+\le2$, irrespective of the much larger
coefficient-saturation index.
\end{corollary}
\begin{proof}
Use the three-term relation
\[
 xL_h=\frac{h+1}{2(2h+1)}L_{h+1}+\frac12L_h
             +\frac h{2(2h+1)}L_{h-1}.
\]
All denominators and the needed nonzero off-diagonal coefficients
are units under the size hypothesis. An interior $xL_h$ adds
nothing. A boundary source supplies exactly its one missing
neighbor; at $h=0$ there is no lower neighbor. Distinct Legendre
degrees in this short high range are independent over
$\mathbb F_p$, and multiplication by $F_1$ is injective.
\end{proof}

If $\deg(FG)<p^2$, the elementary leading-residue identity is
\[
 (\alpha,p^2\beta,p^2\tau,p\sigma)(P)
 \equiv\operatorname{diag}(1,(-1/p),(-3/p),1)
                 C(\mathcal C_pP)\pmod p.                 \tag{\ref{sat:reiwa:classification}.2}
\]
All four rows on the left are integral. This identity follows by
separating multiples of $p$ in the monomial integral recurrence
\[
 J_0=\pi/3,\quad J_1=4-2\sqrt3,\quad
 jJ_j=4(j-1)J_{j-2}-2\sqrt3,
\]
and using $(J_0-J_j)/j$ for the nonconstant monomial conic
integral. Equivalently it follows from the resonant polynomials
for $\mathcal D$; it is the same leading-residue identity used
in Reiwa R49. Consequently Theorem~\ref{sat:reiwa:classification}
gives an exact finite-dimensional test for whether the whole
scaled four-coordinate image is $\mathbb Z_{(p)}^4$: the
listed high coordinate vectors must have rank four modulo $p$.
Surjectivity modulo $p$ lifts by a unit minor, and failure of
rank four prevents surjectivity.

\subsection{Closure of the two-block saturation gap}
Specialize to
\[
 n=100k,\quad K=30k,\quad a=21k,\quad b=k,
 \qquad F=(t-2)^{21k}(t+2)^k.                           \tag{\ref{sat:reiwa:classification}.3}
\]
Keio R51~\cite{K51} proves the following nested-image input: there is a
fixed nonzero integer polynomial $A(k)$ such that, for every
compact interval $I\subset(19/100,39/200)$, all sufficiently
large $k$ and all primes $p/n\in I$ with $p\nmid A(k)$,
\[
 \Gamma(\Lambda)_{(p)}=p^{-1}\mathbb Z_{(p)}
                                      \times p\mathbb Z_{(p)}. \tag{\ref{sat:reiwa:classification}.4}
\]
Its lower inclusion uses the first nine nested columns together
with $FL_{6p+1}$; its endpoint functional is $J_2$, not $J_1$.
The present theorem adds the previously missing saturated upper
inclusion.

\begin{theorem}[Exact saturated equality on the two-block interval]
\label{sat:reiwa:gapclosed}
With the same polynomial exceptions and compact-interval convention
as the nested-image input, one has
\[
 \Gamma(\Lambda^{\rm sat})_{(p)}
 =\Gamma(\Lambda)_{(p)}
 =p^{-1}\mathbb Z_{(p)}\times p\mathbb Z_{(p)}
 \qquad(19/100<p/n<39/200).
\]
In particular saturation neither enlarges the pair image nor
increases its least beta denominator at these primes.
\end{theorem}
\begin{proof}
Here $a=p+s$, $b=w$, $m=22k-p<p/2$, and
$\mathcal H=\{5,6\}$. A nonzero divided source at $h=5$ would
require
\[
 4p+22k-1\ge100k,
\]
which is impossible for $p<(39/2)k$. At $h=6$ it would require
\[
 8p-22k\le130k,
\]
which is impossible for $p>19k$. Thus $\mathcal J$ is empty.
The entire saturated selected image is exactly the span of
$q_5=(x-2)L_5$ and $q_6=(x-2)L_6$.

Their two high nuisance coordinates are independent, since the
canonical triangular recurrence gives the exact determinant
\[
 \det\begin{pmatrix}\tau(q_5)&\sigma(q_5)\\
                     \tau(q_6)&\sigma(q_6)\end{pmatrix}
             =-\frac{150534053}{315}\ne0.
\]
Absorb its fixed numerator and denominator primes into $A$.
For an exactly nuisance-free saturated input, the leading-residue
identity therefore forces $\mathcal C_pP=0$. The alpha row gives
$\alpha\in p\mathbb Z_{(p)}$ and the beta row gives
$\beta\in p^{-1}\mathbb Z_{(p)}$. This is the desired upper
inclusion for \emph{every} saturated input. The nested-image
input supplies the reverse inclusion.
\end{proof}

\begin{corollary}[No exponential gain from the closed gap]
\label{sat:reiwa:globalgap}
Let $\Lambda^{\mathcal I}$ be the partial saturation at all primes
with $19n/100<p<39n/200$. Let $G_n$ and $G_n^{\mathcal I}$ be
the least positive integers clearing the beta coordinates after
division by $2^{22k}$ on the nested and partially saturated
nuisance-free pair images. Then
\[
 \log[\Gamma(\Lambda^{\mathcal I}):\Gamma(\Lambda)]=o(n),
 \qquad \log(G_n^{\mathcal I}/G_n)=o(n).
\]
For either successive minimum and any common, possibly
$n$-dependent, pair norm, the two genuinely normalized lattices
have minimum ratio $\exp(o(n))$.
\end{corollary}
\begin{proof}
The Legendre expansion coefficient of an integer polynomial of
degree at most $d=n+K$ is
$(2j+1)\int_0^1GL_j$. Consequently
\[
 B_d\Lambda^{\rm sat}\subseteq\Lambda,
 \qquad B_d=\operatorname{lcm}(1,\ldots,2d+1).
\]
For proportional primes in question, $p^2>2d+1$ eventually;
the rank-two pair index is thus at most $p^2$ at each prime.
On every fixed compact subinterval, Theorem~\ref{sat:reiwa:gapclosed}
leaves only primes dividing $A(k)$, whose total logarithmic
contribution is $O(\log k)$. The two excluded boundary strips
of relative width $\varepsilon$ contribute at most
$4\varepsilon n+o(n)$ by the prime number theorem. Let
$\varepsilon$ tend to zero after the limit in $n$.

Write the pair index as $J_n$ and the denominator ratio as $q_n$.
Then $q_n\mid J_n$, and for the normalized images $Z,Z'$,
\[
 q_nZ\subseteq Z',\qquad J_nZ'\subseteq q_nZ.
\]
Thus $(q_n/J_n)\lambda_i(Z)\le\lambda_i(Z')\le q_n\lambda_i(Z)$,
which proves the claim for every common norm.
\end{proof}

This is a closed local-image and global-comparison question, not
an irrationality assertion. It removes an unclassified candidate
source of arithmetic improvement. It does not supply a small
nonzero surviving form, nor bound the second minimum itself.


% New synthesis: quantitative arithmetic for the Keio normality theorems.
% Requires amsmath, amssymb, amsthm and theorem/lemma/corollary environments.
\section{Quantitative arithmetic of the conic normalizers}
\label{height:keio:newheight}

The all-index normality theorems remove rational zeros from the high-block
construction, but rational nonvanishing alone does not control the primes
excluded when the high index grows. We close this quantitative gap. The
argument uses only the polynomially specified coordinates, so it makes no
assumption concerning rational independence of the real periods.

Put $B_d=\operatorname{lcm}(1,\ldots,d)$, with $B_0=1$. Let
$\rho(t)=2/\sqrt{4-t^2}$. Define the density coordinates by the unique
polynomial reduction
\[
 P=c+\tfrac12((4-t^2)d_P'-t d_P),\qquad
 \mathcal J(P)=(-2d_P(0),d_P(1),c/3).
\]
Define the conic coordinates by
\[
 P=P(0)+\tfrac t2((4-t^2)U_P''-3tU_P'-U_P),\qquad
 \mathcal I(P)=\left(2U_P(0),P(0),-U_P(1),
                         \frac{3U_P'(1)-U_P(1)}6\right).
\]
We order these coordinates as $(\beta,\alpha,\tau,\sigma)$.
For a reduced rational number $z=a/b$, $b>0$, put
$H(z)=\max(|a|,b)$; in particular $H(0)=1$.

\begin{lemma}[Uniform coefficient arithmetic]\label{height:keio:lcm}
If $P\in\mathbb Z[t]$ has degree at most $d$ and coefficient height at
most $H\ge1$, then
\[
 3B_d\mathcal J(P)\in\mathbb Z^3,\qquad
 3B_d^2\mathcal I(P)\in\mathbb Z^4.
\]
Moreover
\[
 \|\mathcal J(P)\|_\infty,\ \|\mathcal I(P)\|_\infty
       \le 14H2^d.
\]
Consequently the reduced heights of all these coordinates are at most
$42H2^dB_d^2=H\exp(O(d))$, with an absolute implied constant.
\end{lemma}
\begin{proof}
For $j\ge1$ the nonzero coefficients in $d_{t^j}$ are
\[
 [t^{j-1-2r}]d_{t^j}
 =-2\,4^r\,
 \frac{(j-1)(j-3)\cdots(j-2r+1)}
      {j(j-2)\cdots(j-2r)},
 \qquad 0\le r\le\lfloor(j-1)/2\rfloor.
\]
For an odd prime $p$, and for each $p^a\le j$, the number of multiples
of $p^a$ in the denominator progression exceeds their number in the
numerator progression by at most one. Indeed these are progressions
of the same invertible step, of lengths $r+1$ and $r$; each count
lies between the floor and ceiling of its length divided by $p^a$.
Thus the denominator valuation is at most
$\lfloor\log_p j\rfloor$.
If $j$ is odd the displayed denominator is odd. If $j=2m$, its
2-adic valuation is bounded below by
\[
 r-v_2\big(m(m-1)\cdots(m-r)\big)
 =r-v_2((r+1)!)-v_2\binom m{r+1}
 \ge-\lfloor\log_2m\rfloor.
\]
Here $v_2((r+1)!)\le r$, and the factorial-floor formula bounds the
binomial valuation by the number of powers of two at most $m$.
Hence $B_jd_{t^j}$ is integral. The remaining constant is
$c(t^{2r})=\binom{2r}r$ and $c(t^{2r+1})=0$, as follows either from
the same triangular reduction or from the full density integral.
This proves the assertion for $\mathcal J$.

Writing $\mathcal J_j=\mathcal J(t^j)$, direct reduction gives
\[
 \mathcal J_0=(0,0,1/3),\quad \mathcal J_1=(4,-2,0),\quad
 \mathcal J_j=\frac{4(j-1)}j\mathcal J_{j-2}-(0,2/j,0)
 \quad(j\ge2).
\]
The conic rows satisfy $\mathcal I(1)=(0,1,0,0)$ and, for $j>0$,
\[
 \mathcal I(t^j)=\frac1j
       (-\mathcal J_j^\beta,0,-\mathcal J_j^\tau,
                                      1/3-\mathcal J_j^\sigma).
\]
Since $j\mid B_d$, this proves the second denominator assertion.
The density recurrence gives $\|\mathcal J_j\|_\infty\le6\,2^j$
by induction, and the last formula gives
$\|\mathcal I(t^j)\|_\infty\le7\,2^j$. Summing coefficients of
$P$ proves the displayed bounds. Finally $\log B_d=O(d)$, by the
elementary Chebyshev bound for the least common multiple. This bound
also follows by summing the central-binomial bound for prime powers
over dyadic intervals.
\end{proof}

% Optional self-contained core, with the coordinate definitions of keio_new.tex.
\subsection{The normality mechanism}
\label{height:keio:normalityproof}
Here is a proof of the normality statements on which the fixed-window
inversion and high-block arithmetic rest. It also specifies exactly why
formal coordinates can be used without an arithmetic assumption about $V$.

\begin{lemma}[One extra positive-density condition]
\label{height:keio:signchange}
Suppose a nonzero real polynomial $Q$ is orthogonal on $(0,1)$ to
polynomials of degree below $k$, and $\int_0^1Q(t)\rho(t)\,dt=0$.
Then $Q$ has at least $k+1$ sign changes in $(0,1)$.
\end{lemma}
\begin{proof}
For $k=0$ use positivity of $\rho$. If the number of sign changes is
less than $k$, their product polynomial contradicts ordinary
orthogonality. If it is exactly $k$, interpolate $\rho$ at these nodes
by a polynomial $p$ of degree below $k$. The interpolation error has
the sign of the node product, because $\rho^{(k)}(t)>0$ on $(0,1)$.
The latter follows from the positive power series for
$(1-t^2/4)^{-1/2}$. Thus $Q(\rho-p)$ has a fixed nonzero sign apart
from finitely many zeros, contrary to its zero integral.
\end{proof}

\begin{theorem}[Three density columns and four conic columns]
\label{height:keio:normality}
For every $h\ge1$ and every $n\ge0$,
\[
 \det\bigl(\mathcal J(L_{h-1}),\mathcal J(L_h),
                         \mathcal J(L_{h+1})\bigr)\ne0,
\]
\[
 \det\bigl(\mathcal I(L_n),\mathcal I(L_{n+1}),
              \mathcal I(L_{n+2}),\mathcal I(L_{n+3})\bigr)\ne0.
\]
In these formulas the indicated coordinate vectors are the rows.
\end{theorem}
\begin{proof}
The density reduction gives the exact antiderivative
\[
 \rho P=c\rho+(\sqrt{4-t^2}\,d_P)' .
\]
Consequently $\mathcal J(P)=0$ forces separately zero integrals of
$P\rho$ over $(-2,0),(0,1),(1,2)$. If a nonzero combination $Q$ of
$L_{h-1},L_h,L_{h+1}$ had zero density coordinates, its ordinary
orthogonality below degree $h-1$ and Lemma~\ref{height:keio:signchange}
would give $h$ inner sign changes. Each outer interval adds one.
This contradicts $\deg Q\le h+1$.

For the four-column assertion, suppose a nonzero combination $Q$ of
$L_n,\ldots,L_{n+3}$ has $\mathcal I(Q)=0$. Its alpha coordinate
implies $Q(0)=0$. Define
\[
 R(t)=\int_0^t Q(s)/s\,ds,\qquad Z(t)=R(t)-R(1).
\]
The monomial identity in Lemma~\ref{height:keio:lcm}, including its constant
$1/3$ term, gives exactly $\mathcal J(Z)=-\mathcal I_3(Q)=0$,
where $\mathcal I_3=(\beta,\tau,\sigma)$. Also $Q=tZ'$ and
$Z(1)=0$. For $0\le j<n$, integration by parts transfers the ordinary
orthogonality:
\[
 0=\int_0^1t^jQ(t)dt=-(j+1)\int_0^1t^jZ(t)dt.
\]
The density lemma and the two outer intervals give $n+3$ distinct
roots of $Z$ in those open intervals; $Z(1)=0$ gives one more.
This contradicts $\deg Z\le n+3$. Since $Q\ne0$, also $Z\ne0$,
so the contradiction is valid.
\end{proof}

Let $L_h(t)=P_h(2t-1)$ be the shifted Legendre polynomial and set
\[
 D_h=\det\begin{pmatrix}
 \mathcal J(L_{h-1})\\\mathcal J(L_h)\\\mathcal J(L_{h+1})
 \end{pmatrix}\quad(h\ge1),\qquad
 T_h=\tau(tL_h),
\]
\[
 E_h=\det\begin{pmatrix}\beta(L_h)&\beta(tL_h)\\
                         \tau(L_h)&\tau(tL_h)\end{pmatrix},
 \qquad
 N_h=\det\begin{pmatrix}\tau(L_h)&\tau(tL_h)\\
                         \sigma(L_h)&\sigma(tL_h)\end{pmatrix}.
\]
The all-index results of Keio R49--R50~\cite{K49,K50} give
\[
 D_h\ne0,\qquad (-1)^hT_h>0,\qquad
 (-1)^hE_h>0,\qquad(-1)^hN_h>0.
\]
These hypotheses are used as established algebraic normality theorems,
not as conclusions drawn from numerical approximation to the periods.

\begin{theorem}[Uniform high-index heights and exceptional primes]
\label{height:keio:badprime}
There is an absolute constant $C>0$ such that, for every $h\ge1$, the
reduced numerator and denominator heights of $D_h,T_h,E_h,N_h$ are at
most $\exp(Ch)$. In particular each of their absolute values is at
least $\exp(-Ch)$.

Let $a_h,b_h,C_h$ be the unique rational normalizers determined by
\[
 w_h=\sigma+a_h\tau+b_h\beta,\quad
 w_h(L_h)=w_h(tL_h)=0,\qquad C_h=\beta(tL_h)/T_h.
\]
Their heights are at most $\exp(Ch)$ after increasing $C$.
Write $v_h=\beta-C_h\tau$ and
\[
 S_h=\det\begin{pmatrix}
 w_h(\mathcal J(L_h))&w_h(\mathcal J(tL_h))\\
 v_h(\mathcal J(L_h))&v_h(\mathcal J(tL_h))
 \end{pmatrix},
\]
where the covectors act in density coordinate order $(\beta,\tau,\sigma)$.
Then $S_h=D_h/(4(2h+1)T_h)$ is nonzero and has the same exponential
height bound.

There is consequently a positive integer $\mathfrak B_H$ with
\[
 \log\mathfrak B_H=O(H^2)
\]
such that for every $1\le h\le H$ and every prime
$p\nmid\mathfrak B_H$, all these normalizers are $p$-integral and
$D_h,T_h,E_h,N_h,S_h$ are $p$-units.
\end{theorem}
\begin{proof}
The coefficient identity
\[
 [t^j]L_h=(-1)^{h-j}\binom hj\binom{h+j}j
\]
gives coefficient height at most $8^h$. Lemma~\ref{height:keio:lcm}, with
degree at most $h+1$, therefore bounds every coordinate needed here
by reduced height $\exp(O(h))$. Determinants of fixed size and a
fixed number of rational field operations preserve such a bound.
Cramer's rule applies because $E_h\ne0$, and $C_h$ is defined because
$T_h\ne0$. The identity for $S_h$ is the exact contact determinant
identity of R49, followed by the two normalizing row operations;
it also shows its nonvanishing. A nonzero reduced rational of height
at most $\exp(Ch)$ has absolute value at least $\exp(-Ch)$.

For each $h$, multiply the nonzero numerators and denominators of
$D_h,T_h,E_h,N_h,S_h$, and the denominators of $a_h,b_h,C_h$, taking
absolute values. The product is a positive integer of logarithm
$O(h)$. Let $\mathfrak B_H$ be the product of these integers for
$1\le h\le H$. Its logarithm is $O(\sum_{h\le H}h)=O(H^2)$,
and its prime-divisor property is exactly the asserted one.
\end{proof}

\begin{corollary}[Four-column rational inversion]
\label{height:keio:fourinverse}
Let $A_h$ have rows $\mathcal I(L_h),\ldots,\mathcal I(L_{h+3})$.
The all-index four-column normality theorem implies
\[
 \max_{i,j} H((A_h^{-1})_{ij})\le\exp(O(h+1)).
\]
Thus four consecutive columns realize any prescribed rational formal
coordinate vector, with a rational control of height at most the input
height to an absolute fixed power times $\exp(O(h+1))$.
\end{corollary}
\begin{proof}
Every entry of $A_h$ has height $\exp(O(h+1))$ by
Lemma~\ref{height:keio:lcm}. Its nonzero determinant and its cofactors use
only a fixed number of rational arithmetic operations. The adjugate
formula proves the bound. Multiplication by the prescribed vector is
again a fixed number of rational operations.
\end{proof}

\begin{corollary}[A growing-block exclusion budget]\label{height:keio:growing}
If $H=H(n)=o(\sqrt n)$, all high-index normalizer exceptions for
$1\le h\le H$ can be removed at total prime logarithmic weight
$o(n)$. More explicitly,
\[
 \sum_{p\mid\mathfrak B_H}\log p=O(H^2),\qquad
 \#\{p\ge Y:p\mid\mathfrak B_H\}
                  =O(H^2/\log Y)\quad(Y\ge2).
\]
Thus these particular exceptional primes do not alter a denominator
exponent measured on the scale $n$ in that growing-block range.
\end{corollary}
\begin{proof}
The sum is bounded by $\log\mathfrak B_H$. Each term with $p\ge Y$
is at least $\log Y$. Apply Theorem~\ref{height:keio:badprime}.
\end{proof}

This conclusion controls high rational normalizers only. It supplies
neither the low-image rank needed for every new prime interval nor a
short integral surviving form. In particular it cannot replace the
second-successive-minimum estimate required for irrationality.


% Self-contained section fragment. All labels have prefix hs:.
\section{An exact dyadic index and an obstruction to short integral output bases}
\label{odd:hs:sec}

The direct odd output basis of the later Heisei manuscripts has exponential
coefficient denominators. A different invariant measures the simultaneous cost
of making an entire output basis integral. We compute its dyadic part exactly.
This is a new consequence of the direct boundary-jet construction, rather than
a claim about small forms containing the Gieseking volume.

Put
\[
 \mathcal B(f,g)=\int_0^1\!\int_0^1
       \frac{f(s)g(t)}{\sqrt{1-s^2t^2/4}}\,ds\,dt,
 \qquad N=n-1,\qquad D_n(z)=\prod_{h=0}^{n-1}(z-2h).
\]
For $0\le i<n$, define
\begin{align}
 f_i(x)&=\frac{1}{(-2)^N(N-i)!}
   \sum_{h=0}^i\frac{(-1)^{i-h}(x-4)^h}
   {(2h)!(N-h)!(i-h)!},\label{odd:hs:f}\\
 w(x)&=(4-x)^{1/2}(1-x)^{2n},\qquad
 Z_i(x)=[w^{-1}f_i]_{<n},\qquad
 P_i(x)=(1-x)^{2n}Z_i(x),\label{odd:hs:Z}\\
 U_nP&=(4-x)^{n-1/2}\frac{d^n}{dx^n}\big((4-x)^{1/2}P\big),\\
 G_{n,i}(t)&=D_n(t\partial_t)\big[t(U_nP_i)(t^2)\big].\label{odd:hs:G}
\end{align}
The notation $[\ ]_{<n}$ denotes Taylor truncation at zero. The square-root
branch has constant term $2$. All polynomials $G_{n,i}$ have rational
coefficients, are odd, and have degree at most $6n-1$.
In the coordinate $y=4-x$, the apparently algebraic operator is diagonal:
\begin{equation}\label{odd:hs:Udiag}
 U_n(y^r)=(-1)^n(r+\tfrac12)_{\underline n}y^r.
\end{equation}
Let $C_n$ be their $(3n)\times n$ coefficient matrix, with rows indexed by
$t,t^3,\ldots,t^{6n-1}$, and let
\[
 \mathscr W_n=\operatorname{span}_{\mathbb Q}
          \{G_{n,0},\ldots,G_{n,n-1}\}.
\]

\begin{lemma}[The output identities]\label{odd:hs:output}
For $0\le h,i<n$,
\[
 \mathcal B(s^{2h},G_{n,i})=\delta_{hi}.
\]
These values are rational by formal cancellation of the $\pi$ and $\sqrt3$
coordinates; no linear independence involving the volume is assumed.
\end{lemma}
\begin{proof}
The explicit adjoint polynomial for row $i$ is
\[
 h_i(x)=(-2)^Ni!(N-i)!\sum_{h=i}^N
             \binom{2h}{h}\binom hi(4-x)^{N-h}.
\]
For any polynomial $Z$ of degree less than $n$, set $F=wZ$ and
$P=(1-x)^{2n}Z$. Direct integration of the odd monomial moments, followed
by the adjoint identity for $D_n(z)/(z-2i)$, gives
\[
 \mathcal B\bigl(s^{2i},D_n(t\partial_t)[t(U_nP)(t^2)]\bigr)
       =-\int_0^1 h_i(x)F^{(n)}(x)\,dx.
\]
Here the endpoint at $1$ vanishes to order $2n$. Integrating by parts $n$
times gives
\[
 \sum_{k=0}^N(-1)^kh_i^{(k)}(0)F^{(N-k)}(0).
\]
There is no integral remainder because $\deg h_i<n$. In terms of
$f=[F]_{<n}$ this boundary expression is
\begin{equation}\label{odd:hs:T}
 (Tf)_i=(-2)^Ni!(N-i)!\sum_{h=i}^N
      \binom{2h}{h}\binom hi(N-h)!f^{(h)}(4).
\end{equation}
Upper triangular binomial inversion shows directly that the polynomials
in \eqref{odd:hs:f} satisfy $Tf_i=e_i$. Since $[wZ_i]_{<n}=f_i$, this proves
the asserted identities. The same calculation is the direct boundary-jet
proof in Heisei R48~\cite{H48}; in particular it evaluates the actual integral.
The formal odd moment expansion belongs to
$\mathbb Q+\mathbb Q\pi+\mathbb Q\sqrt3$. Since
$1,\pi,\sqrt3$ are linearly independent over $\mathbb Q$, its two
nuisance coordinates vanish. Only this classical independence is used;
the volume does not occur in these odd moments.
\end{proof}

\begin{theorem}[Exact dyadic integral-output index]\label{odd:hs:binary}
Let $s_2(j)$ be the number of ones in the binary expansion of $j$, with
$s_2(0)=0$. The lattice
\[
 \Gamma_{n,2}=\{v\in\mathbb Z_2^n:C_nv\in\mathbb Z_2^{3n}\}
\]
has exact index
\begin{equation}\label{odd:hs:index2}
 [\mathbb Z_2^n:\Gamma_{n,2}]
       =2^{A_n},\qquad
 A_n=4n^2-2n-3\sum_{j=0}^{n-1}s_2(j).
\end{equation}
Equivalently, the total negative Smith valuation of the coefficient map
is $A_n$. All its Smith valuations are nonpositive. This does not assert
an individual formula for every Smith factor.
\end{theorem}
\begin{proof}
We first compute the sum of its Smith valuations. The operator
$D_n(t\partial_t)$ multiplies each odd monomial by a product of odd
integers, so it is unimodular over $\mathbb Z_2$. By \eqref{odd:hs:Udiag},
$2^nU_n$ is unimodular in the $y$ monomial basis. The coordinate change
$y=4-x$ and its inverse are integral, so the same is true in the $x$ basis.
Multiplication by $(1-x)^{2n}$ from degree less than $n$ to degree less
than $3n$ is a saturated integral embedding: its first $n$ coefficient
rows form a unit lower triangular matrix. Thus its matrix is equivalent
over $\mathbb Z_2$ to a column inclusion.

Let $F_n$ denote the square coefficient matrix of $f_0,\ldots,f_{n-1}$,
and $B_n$ the matrix of truncated multiplication by $w^{-1}$ on
polynomials of degree less than $n$. The preceding unimodular operations
reduce the nonzero Smith valuations of $C_n$ to those of
\[
 M_n=2^{-n}B_nF_n.
\]
The matrix $B_n$ is lower triangular with diagonal $1/2$, hence
$v_2(\det B_n)=-n$. In the $(x-4)$ monomial basis, $F_n$ is upper
triangular with diagonal
\[
 \frac{1}{(-2)^N(N-i)!^2(2i)!},\qquad 0\le i<n.
\]
The change to the $x$ basis has determinant one. Consequently
\begin{align*}
 v_2(\det M_n)
 &=-n^2-n-nN-2\sum_{i=0}^{n-1}v_2(i!)
                         -\sum_{i=0}^{n-1}v_2((2i)!)\\
 &=-4n^2+2n+3\sum_{i=0}^{n-1}s_2(i)=-A_n,
\end{align*}
by $v_2(i!)=i-s_2(i)$ and $s_2(2i)=s_2(i)$.

A determinant alone would not identify the integral-preimage index if
positive Smith valuations were present. We exclude that possibility
using the actual rational-coordinate map. For the odd monomial
$t^{2r+1}$ its rational coordinate on row $s^{2h}$ is
\begin{equation}\label{odd:hs:Rentry}
 R_{h,r}=\frac{4\,16^r(r!)^2}
                   {(2r+1)!(2h-2r-1)}.
\end{equation}
For clarity, this follows by inserting
$d_{2r+1}(0)=-2\,16^r(r!)^2/(2r+1)!$ into the formal moment identity
\[
 \mathcal B(s^{2h},t^j)=\frac{J_j-J_{2h}}{2h-j},\qquad
 J_j=\frac\pi3c_j+\sqrt3d_j(1)-2d_j(0),
\]
where $j$ is odd and $d_{2h}(0)=0$. Every entry in \eqref{odd:hs:Rentry}
is $2$-integral, since
\[
 v_2(R_{h,r})=4r+2-s_2(r)\ge0.
\]
Lemma~\ref{odd:hs:output} and its formal cancellation give $R_nC_n=I_n$.
Thus $C_n$ has an integral left inverse. In its Smith normal form a
positive exponent would require a nonintegral entry in any left inverse.
All exponents are therefore nonpositive. The index of the integral
preimage is the product of $2$ to their negatives, which is
$2^{-v_2(\det M_n)}=2^{A_n}$.
\end{proof}

The theorem concerns the output lattice, not merely a common denominator
of one chosen basis. In particular it remains true after an integral
unimodular change of output coordinates. Since
$\sum_{j<n}s_2(j)=O(n\log(n+1))$, its logarithm is
$4(\log2)n^2+O(n\log(n+1))$.

\paragraph{Combining the dyadic index with the odd-prime residue ranks.}
The sharp residue-rank theorem of Heisei R49~\cite{H49} supplies an independent
contribution. We state its precise input here. For an odd prime $p>n$,
all coefficient poles are at most simple. If
$p\le2n-2$, write $p=2q+1$ and $n=q+1+a$. Then
\begin{equation}\label{odd:hs:odd-rank}
 r_{n,p}:=\operatorname{rank}_{\mathbb F_p}(pC_n\bmod p)
       =\min\{a,\max(0,q-2a-1)\};
\end{equation}
outside that range $r_{n,p}=0$. Thus the local integral-preimage index is
$p^{r_{n,p}}$. This follows from the Taylor-tail map
$r\mapsto(1-x)^dr-[(1-x)^dr]_{<n}$ on $\deg r<a$, where
$d=p-1-2a$. Its rank is the right side of \eqref{odd:hs:odd-rank}; the
remaining two operators are injective on its tail image by the three
root-multiplicity argument at $0,1,4$. No determinant-to-rank inference
is involved.

Define the global lattice
\[
 \Gamma_n=\{v\in\mathbb Z^n:C_nv\in\mathbb Z^{3n}\}.
\]
It has finite index because $C_n$ is rational. Localizing and using
Theorem~\ref{odd:hs:binary} and \eqref{odd:hs:odd-rank} gives
\begin{equation}\label{odd:hs:global-lower}
 \log[\mathbb Z^n:\Gamma_n]
 \ge\left(4\log2+\frac1{12}\right)n^2+o(n^2).
\end{equation}
Indeed, uniformly for $p>n$,
\[
 \frac{r_{n,p}}n=\psi(p/n)+O(1/n),\qquad
 \psi(t)=
 \begin{cases}
  \frac32t-2,&4/3\le t\le3/2,\\
  1-\frac12t,&3/2\le t\le2,\\
  0,&\text{otherwise}.
 \end{cases}
\]
The prime number theorem and partial summation give
$\sum_{p>n}r_{n,p}\log p=n^2\int\psi(t)dt+o(n^2)
=n^2/12+o(n^2)$. No assertion about the exact contribution of the
remaining odd primes is needed for \eqref{odd:hs:global-lower}.

\begin{corollary}[No uniformly subexponential integral output basis]
\label{odd:hs:height-obstruction}
Suppose $Q_{n,1},\ldots,Q_{n,n}\in\mathscr W_n\cap\mathbb Z[t]$
are linearly independent and all their rational output vectors
\[
 v_j=(\mathcal B(1,Q_{n,j}),\mathcal B(s^2,Q_{n,j}),\ldots,
                   \mathcal B(s^{2n-2},Q_{n,j}))
\]
belong to $\mathbb Z^n$. For the coefficient norm
$\|Q\|_1=\sum_k|[t^k]Q|$,
\begin{equation}\label{odd:hs:height}
 \left(\prod_{j=1}^n\|Q_{n,j}\|_1\right)^{1/n}
 \ge \exp\left(\left(4\log2+\frac1{12}+o(1)\right)n\right).
\end{equation}
In particular at least one control satisfies the same lower bound.
The constant is $4\log2+1/12=2.8559220555\ldots$.
\end{corollary}
\begin{proof}
By Lemma~\ref{odd:hs:output}, the $v_j$ are the coordinates of the $Q_{n,j}$
in the direct output basis. They are independent elements of $\Gamma_n$.
Therefore their determinant has absolute value at least
$[\mathbb Z^n:\Gamma_n]$. The integral kernel is at most $2/\sqrt3$
on the unit square, so
\[
 |\mathcal B(s^{2h},Q)|\le\frac{2\|Q\|_1}{\sqrt3(2h+1)},
 \qquad
 \|v_j\|_2\le\frac\pi{\sqrt6}\|Q_{n,j}\|_1,
\]
using $\sum_{h\ge0}(2h+1)^{-2}=\pi^2/8$. Hadamard's inequality
now gives
\[
 [\mathbb Z^n:\Gamma_n]\le
       \left(\frac\pi{\sqrt6}\right)^n
                  \prod_{j=1}^n\|Q_{n,j}\|_1.
\]
Insert \eqref{odd:hs:global-lower} and take $n$th roots.
\end{proof}

This obstruction has a precise scope. It concerns a full set of $n$
independent controls in the specified odd space, with both polynomial
coefficients and rational outputs integral. It does not bound every
single control, does not include arbitrary larger windows, and does not
exclude small target-bearing forms. It shows that the exponential cost
of a complete integral output basis cannot be removed merely by choosing
a better rational basis. In contrast, a single small surviving form is
exactly the arithmetic task that remains relevant to irrationality.



\section{A calibrated one-row moment lattice}
\subsection{Moment coordinates and the assertion}
Put
\[
 \B(F,G)=\int_0^1\!\int_0^1
 \frac{F(s)G(t)}{\sqrt{1-s^2t^2/4}}\,ds\,dt,
 \qquad V=\B(1,1)=\operatorname{Cl}_2(\pi/3).
\]
The last equality follows by integrating in $s$, then putting
$t=2\sin(u/2)$ and integrating by parts. In particular $1<V<6/5$.
Write $J_j=\int_0^1t^j(1-t^2/4)^{-1/2}\,dt$ and
$I_j=\int_0^1-t^j\log t\,(1-t^2/4)^{-1/2}\,dt$.
The formal coordinate order is $(1,\sqrt3,\pi,V)$, with
\begin{align}
 J_0&=\pi/3,&J_1&=4-2\sqrt3,\notag\\
 J_j&=\frac{4(j-1)}jJ_{j-2}-\frac2j\sqrt3,\label{calbase:Jrec}\\
 I_0&=V,&I_j&=\frac{4(j-1)}jI_{j-2}
              -\frac4jJ_{j-2}+\frac1jJ_j\quad(j\ge2\text{ even}).\label{calbase:Irec}
\end{align}
These identities follow by integration by parts; differentiating the
exponent in the first recurrence proves the second. Moreover
\[
 \B(s^i,t^j)=\frac{J_j-J_i}{i-j}\ (i\ne j),\qquad
 \B(s^i,t^i)=I_i.
\]
All statements about coordinates below concern this specified formal
expansion, not an assumed linear independence of the four constants.

Fix $a\ge1$, put $m=2a$ and $C=\binom{2a}{a}$. Let
$\mathcal P_m=t^m\mathbb Q[t]_{\le3}$.
Let $L_k(t)=P_k(2t-1)$ be the shifted Legendre polynomial. For
$q=\sum_k u_kL_k$, define the rational norm
\begin{equation}
 \N_m(q)=\frac{11}{10(m+1)}\sum_k |u_k|(5/29)^k.
 \label{calbase:norm}
\end{equation}

\begin{theorem}\label{calbase:main}
The four-coordinate map $q\mapsto\B(s^m,q)$ is an isomorphism
$\mathcal P_m\longrightarrow\mathbb Q^4$.
Let $q_\beta,q_\alpha\in\mathcal P_m$ have formal outputs $1,V$,
respectively, and set
$\Lambda_m=\mathbb Zq_\beta+\mathbb Zq_\alpha$.
For every real $b,d$,
\begin{equation}
 |b+Vd|\ \le\ \N_m(bq_\beta+dq_\alpha)
 \ \le\ 32|b+Vd|+\frac{104|d|}{C(m+1)^2}.
 \label{calbase:calibration}
\end{equation}
If $\lambda_2(m)$ is the second successive minimum of $\Lambda_m$
for this norm, then
\[
 \lambda_2(2a)\longrightarrow0
 \quad\Longleftrightarrow\quad V\notin\mathbb Q.
\]
If $V=P/Q$ in lowest terms, $Q>0$, then
$\lambda_2(2a)\ge1/Q$ at every size.
\end{theorem}

\subsection{A rational majorant independent of the target value}
\begin{lemma}\label{calbase:majorant}
For every integer $m\ge0$ and $k\ge0$,
\[
 0<\B(s^m,L_k)\le\frac{\log3}{m+1}(3+2\sqrt2)^{-k}
 <\frac{11}{10(m+1)}(5/29)^k.
\]
Consequently $|\B(s^m,q)|\le\N_m(q)$. If
$q=t^m\sum_{j=0}^3c_jt^j$ and $H=\max_j|c_j|$, then
\begin{equation}\label{calbase:coefnorm}
 \N_m(q)\le\frac{8H}{(m+1)^2}.
\end{equation}
\end{lemma}
\begin{proof}
Rodrigues' formula, with its boundary terms zero, gives
\[
 \int_0^1t^rL_k(t)\,dt=
 \begin{cases}(r!)^2/((r-k)!(r+k+1)!),&r\ge k,\\0,&r<k.
 \end{cases}
\]
Expand the positive kernel in even powers of $st$. Since
$\binom{2r}r/16^r\le4^{-r}$ and $(m+2r+1)^{-1}\le(m+1)^{-1}$,
termwise comparison of these nonnegative Mellin moments bounds
$\B(s^m,L_k)$ by $(m+1)^{-1}\int_0^1L_k(t)/(1-t^2/4)\,dt$.
For $z>1$, another integration by parts gives
\[
 E_k(z):=\int_0^1\frac{L_k(t)}{z-t}\,dt
 =\int_0^1\left(\frac{t(1-t)}{z-t}\right)^k\frac{dt}{z-t}.
\]
The maximum of the parenthesis is
$(\sqrt z-\sqrt{z-1})^2$. Partial fractions and $t\mapsto1-t$
give
\[
 \int_0^1\frac{L_k(t)}{1-t^2/4}\,dt
 =E_k(2)+(-1)^kE_k(3)
 \le\log3\,(3+2\sqrt2)^{-k}.
\]
Here $29/5<3+2\sqrt2$, and $\log3<11/10$ because the
degree-five Taylor sum of $e^{11/10}$ exceeds $3$ by
$15101/12000000$.

Finally $u_k=(2k+1)\int qL_k$ and $|L_k|\le1$ on $[0,1]$,
so $|u_k|\le4H(2k+1)/(m+1)$. Since
\[
 \sum_{k\ge0}(2k+1)(5/29)^k=493/288,
 \qquad (11/10)\cdot4\cdot493/288<8,
\]
equation \eqref{calbase:coefnorm} follows.
\end{proof}

\subsection{An explicit positive calibration}
The recurrences give
\[
 [\pi]J_m=C/3,\quad [1]J_{m+1}=A:=
 \frac{4\cdot16^a(a!)^2}{(2a+1)!},\quad
 [V]I_m=C,\quad [1]I_m=0,
 \quad [\pi]I_m=-Ch/3,
\]
where $h=H_{2a}-H_a\in(0,1)$. Set
\[
 r=\frac{4(m+1)}{m+2},\quad \eta=\frac{4(m+2)}{m+3},\quad
 P(t)=\frac{\eta}{\eta-1}t+\frac2{r-1}t^2-\frac3{\eta-1}t^3,
\]
\[
 Z(t)=\frac{-t+3t^3}{\eta-1},\qquad
 W(t)=1-\frac{2h}{r-1}t^2.
\]
Let $U=\B(s^m,t^mP)$, $T=\B(s^m,t^mZ)$ and
$E=\B(s^m,t^mW/C)$. 
Direct coordinate substitution shows that $t^mP$ has only a
$\sqrt3$ coordinate; $t^mZ$ has rational coordinate $-A$, zero
$\pi,V$ coordinates; and $t^mW/C$ has $V$ coordinate one and
zero rational and $\pi$ coordinates.

For clarity, this calculation can also be done without guessing the
polynomials. If $q=t^m(c_0+c_1t+c_2t^2+c_3t^3)$ has target
coefficient $d$, then $c_0=d/C$. Write
$u=c_1+c_3/3$ and $v=c_1+\eta c_3/3$. Its rational coordinate
is $-Av$, and its $\pi$ coordinate vanishes exactly when
$c_2=2(u-dh/C)/(r-1)$. Thus, after these three coordinates
are fixed, the remaining freedom is exactly addition of $t^mP$.

The quadratic $P(t)/t$ is concave on $[0,1]$. Its endpoint
values and $3\le r<4$, $16/5\le\eta<4$ imply
\begin{equation}\label{calbase:bounds1}
 \frac23t\le P(t)\le\frac52t,\qquad |Z(t)|\le2t,
 \qquad |W(t)|\le2.
\end{equation}
In particular $U>0$. Its formal $\sqrt3$ coordinate is therefore
nonzero, proving the asserted four-coordinate isomorphism.
Writing $v=-1/A$, the unique rational polynomials with outputs
$1,V$ satisfy the exact real identities
\begin{align}
 q_\beta&=t^m\left(\frac{1-vT}{U}P+vZ\right),\label{calbase:qbeta}\\
 q_\alpha&=t^m\left(\frac{V-E}{U}P+\frac WC\right).
 \label{calbase:qalpha}
\end{align}
Although these formulas display real constants, their coefficients are
rational by the coordinate calculation just given. No rationality of
$V$ has been assumed.

The kernel lies between $1$ and $6/5$. Therefore
\begin{equation}\label{calbase:bounds2}
 \frac2{3(m+1)(m+2)}\le U\le\frac3{(m+1)(m+2)},\quad
 |T|\le\frac3{(m+1)(m+2)},\quad
 |E|\le\frac3{C(m+1)^2}.
\end{equation}
Also $A\ge4C$: the sequence $(2a+1)C^2/16^a$ starts at one
and its consecutive ratio is $(2a+1)(2a+3)/(2a+2)^2<1$.
The maximum coefficient heights of $P,Z,W$ are at most
$3/2,3/2,1$, respectively. Equations \eqref{calbase:bounds2} yield
\[
 |T/U|\le9/2,\quad |E/U|\le6/C,
 \qquad H(t^{-m}q_\beta)\le4(m+1)^2.
\]
Subtracting $V$ times \eqref{calbase:qbeta} from \eqref{calbase:qalpha} gives
\[
 q_\alpha-Vq_\beta
 =t^m\left\{\frac WC-\frac EU P
             +Vv\left(\frac TU P-Z\right)\right\}.
\]
Its coefficient height is at most
\[
 \frac1C+\frac9C+
 \frac{6}{5}\frac1{4C}\left(\frac{27}4+\frac32\right)
 =\frac{499}{40C}<\frac{13}C.
\]
Applying Lemma \ref{calbase:majorant} to these two polynomials proves
\eqref{calbase:calibration}, including its lower bound from their actual outputs.

\subsection{The successive-minimum dichotomy}
If $V=P/Q$, two independent integer pairs $(b,d)$ cannot both
lie in the one-dimensional line $Qb+Pd=0$. One of their actual
forms has absolute value at least $1/Q$, and hence so does one
of their norms.

Conversely, suppose $V$ is irrational. For every $\varepsilon>0$
there exist two independent integer pairs with
$0<|b+Vd|<\varepsilon$. This is the elementary pigeonhole
approximation argument: choose a primitive pair of sufficiently
small nonzero error, and then a pair with still smaller nonzero
error, which cannot be an integer multiple of the primitive pair.
Fix these two pairs. Since $C(m+1)^2\to\infty$, their norms
have limsup at most $32\varepsilon$ by \eqref{calbase:calibration}.
Letting $\varepsilon\downarrow0$ proves the converse and the theorem.



% Self-contained extension of the calibrated one-row model, with polynomial norm comparison.
% Parent: use theorem/lemma/corollary environments and amsmath/amssymb.
\section{The exact Diophantine exponent of the calibrated moment lattice}
\label{cal:sec:calibrated-exponent}

This section strengthens the qualitative calibration dichotomy of the source
Meiji R59~\cite{M59}. It determines exactly which Diophantine exponent is measured by
that lattice. The result does not assume that the Gieseking constant is
irrational; the rational alternative is included separately.

Put $V=\operatorname{Cl}_2(\pi/3)$, $m=2a\ge2$, and
$C_m=\binom{2a}{a}$. Let $L_k(t)=P_k(2t-1)$ and write
$q=\sum u_kL_k$. On $t^m\mathbb R[t]_{\le3}$ use the rationally defined norm
\[
 \mathcal N_m(q)=\frac{11}{10(m+1)}
                   \sum_k |u_k|(5/29)^k.
\]
The four formal-coordinate map for
$\mathcal B(s^m,q)$ is invertible on this space. Denote its rational
inverse columns for outputs $1,V$ by $q_{\beta,m},q_{\alpha,m}$.
The established calibration is
\begin{equation}\label{cal:eq:cal-existing}
 |b+Vd|\le\mathcal N_m(bq_{\beta,m}+dq_{\alpha,m})
 \le32|b+Vd|+\frac{104|d|}{C_m(m+1)^2}.
\end{equation}
Moreover the coefficient of $t^m$ of this polynomial is $d/C_m$:
only the diagonal mixed moment has a target coordinate, equal to $C_m$.
These are the only features of the moment construction used below.
Let $\lambda_1(m)\le\lambda_2(m)$ be the successive minima of
$\mathbb Zq_{\beta,m}+\mathbb Zq_{\alpha,m}$ for $\mathcal N_m$.

For an irrational real number $\xi$, its irrationality exponent is
\[
 \mu(\xi)=\sup\{\nu:|\xi-p/q|<q^{-\nu}
       \text{ for infinitely many integers }p,q,\ q\ge1\}.
\]
The value $+\infty$ is permitted, and $1/\infty=0$.

\begin{theorem}[Quantitative calibration]\label{cal:thm:quant-calibration}
For every even $m\ge2$ and every real $b,d$, put
$S_m(b,d)=\max\{|b+Vd|,|d|/C_m\}$. Then
\begin{equation}\label{cal:eq:poly-equivalence}
 \frac{S_m(b,d)}{32(m+7)^{17}}
 \le \mathcal N_m(bq_{\beta,m}+dq_{\alpha,m})
 \le136S_m(b,d).
\end{equation}
If $V$ is irrational, then
\begin{align}
 \limsup_{a\to\infty}
 \frac{-\log\lambda_1(2a)}{\log C_{2a}}
 &=1-\frac1{\mu(V)},\label{cal:eq:lambda1-mu}\\
 \liminf_{a\to\infty}
 \frac{-\log\lambda_2(2a)}{\log C_{2a}}
 &=\frac1{\mu(V)}.\label{cal:eq:lambda2-mu}
\end{align}
If instead $V=P/Q$ in lowest terms, $Q>0$, the two ratios
have limits $1$ and $0$, respectively, and
$\lambda_2(2a)\ge1/Q$ at every size.
\end{theorem}

\begin{proof}[The polynomial norm comparison]
Write $q=t^m\sum_{j=0}^3c_jt^j$ and
$v_k=\int_0^1q(t)t^k\,dt$ for $0\le k\le3$. Then
\[
 v_k=\sum_{j=0}^3\frac{c_j}{m+k+j+1}.
\]
The $4$ by $4$ Cauchy matrix $A_m$ in this identity has determinant
\[
 \det A_m=\frac{144}{\prod_{i,j=0}^3(m+i+j+1)}.
\]
Every cofactor has absolute value at most $6$, by its six-term
expansion and the fact that all entries have absolute value at most one.
Hence every inverse entry is at most $(m+7)^{16}/24$.
For $k\le3$, the expansion of $t^k$ in $L_0,\ldots,L_k$ has
nonnegative coefficients with sum one. Orthogonality therefore gives
$|v_k|\le\sum_{\ell=0}^3|u_\ell|$. Consequently
\[
 \max_j|c_j|\le\frac{(m+7)^{16}}6\sum_{\ell=0}^3|u_\ell|
 \le\frac{10}{66}(29/5)^3(m+7)^{17}\mathcal N_m(q)
 <32(m+7)^{17}\mathcal N_m(q).
\]
Apply this to $c_0=d/C_m$ and combine with the lower bound in
\eqref{cal:eq:cal-existing}. The upper bound follows from the same
calibration by replacing $(m+1)^{-2}$ by one.
\end{proof}

\begin{lemma}[Two-dimensional minima product]\label{cal:lem:two-minima}
For any real $\xi$ and $0<\varepsilon\le1$, let
\[
 S_\varepsilon(b,d)=\max\{|b+\xi d|,\varepsilon|d|\},
 \qquad (b,d)\in\mathbb Z^2,
\]
and let $s_1\le s_2$ be its two successive minima. Then
\[
 \varepsilon/2\le s_1s_2\le2\varepsilon.
\]
\end{lemma}
\begin{proof}
The image lattice
$\{(b+\xi d,\varepsilon d):(b,d)\in\mathbb Z^2\}$
has covolume $\varepsilon$. Two independent vectors of sup norms
$s_1,s_2$ have determinant at least $\varepsilon$ and at most
$2s_1s_2$, proving the lower bound.

For completeness, $s_1\le\sqrt\varepsilon$. Put
$Q=\lfloor\varepsilon^{-1/2}\rfloor$. Among the $Q+1$ points
$0,\xi,\ldots,Q\xi$ modulo one, a circular gap is at most
$1/(Q+1)$. This supplies $1\le d\le Q$ with
$\|d\xi\|\le1/(Q+1)\le\sqrt\varepsilon$, and
$\varepsilon d\le\sqrt\varepsilon$.
Choose a shortest lattice vector $v$, so $\|v\|_\infty=s_1$.
It is primitive and extends to a lattice basis $(v,w)$.
After subtracting an integer multiple of $v$, the parallel component
of $w$ has length at most $\|v\|_2/2$; its perpendicular component
has length $\varepsilon/\|v\|_2$. Therefore
\[
 \|w\|_\infty\le\|w\|_2
 \le s_1/\sqrt2+\varepsilon/s_1
 <2\varepsilon/s_1,
\]
where $s_1^2\le\varepsilon$ was used. The independent pair $(v,w)$
proves the upper bound.
\end{proof}

\begin{proof}[The exponent identities]
Write $\delta(q)=\|q\xi\|$ for irrational $\xi$ and put
$\omega=\mu(\xi)-1$. Equivalently,
\[
 \omega=\limsup_{q\to\infty}\frac{-\log\delta(q)}{\log q}.
\]
The first minimum satisfies
\[
 s_1(\varepsilon)=\min\left(1,\min_{q\ge1}
                         \max\{\delta(q),\varepsilon q\}\right).
\]
If $\omega<\infty$ and $\eta>\omega$, then
$\delta(q)\ge q^{-\eta}$ for all sufficiently large $q$.
The minimum of $\max(q^{-\eta},\varepsilon q)$ over positive
real $q$ is $\varepsilon^{\eta/(\eta+1)}$. The finitely many
remaining $q$ have fixed nonzero errors. Thus
\[
 \limsup_{\varepsilon\downarrow0}
 \frac{-\log s_1(\varepsilon)}{-\log\varepsilon}
 \le\frac\omega{1+\omega}.
\]
For every $0<\eta<\omega$, infinitely many $q$ satisfy
$\delta(q)<q^{-\eta}$. At $\varepsilon_q=\delta(q)/q$,
$s_1(\varepsilon_q)\le\delta(q)$, and
\[
 \frac{-\log s_1(\varepsilon_q)}{-\log\varepsilon_q}
 \ge\frac{-\log\delta(q)}{\log q-\log\delta(q)}
 >\frac\eta{1+\eta}.
\]
Letting $\eta\uparrow\omega$ gives equality. If $\omega=\infty$,
the same lower argument tends to one; the reverse inequality
follows from $s_1(\varepsilon)\ge\varepsilon$ for
$0<\varepsilon\le1$. Thus the first identity holds also for
$\mu(\xi)=\infty$.

These limsup statements remain true on
$\varepsilon_a=C_{2a}^{-1}$. Indeed
\[
 \frac{C_{2a+2}}{C_{2a}}=4-\frac2{a+1}\in[3,4),
\]
so any sufficiently small $\varepsilon$ lies between two consecutive
$\varepsilon_a$ with ratio at most four. Multiplying $\varepsilon$
by a factor in $[1/4,4]$ changes $S_\varepsilon$, and hence its
first minimum, by at most the same factor. The resulting logarithmic
errors are bounded. Lemma~\ref{cal:lem:two-minima} gives
\[
 \frac{-\log s_2(\varepsilon_a)}{\log C_{2a}}
 =1-\frac{-\log s_1(\varepsilon_a)}{\log C_{2a}}+o(1).
\]
Finally \eqref{cal:eq:poly-equivalence} transfers both minima to
$\lambda_i(2a)$ with logarithmic error $O(\log a)$, whereas
$\log C_{2a}=a\log4+O(\log a)$. This proves
\eqref{cal:eq:lambda1-mu}--\eqref{cal:eq:lambda2-mu}.

If $\xi=P/Q$, the primitive zero direction is $(-P,Q)$.
For sufficiently small $\varepsilon$, one has
$s_1=\varepsilon Q$ and $s_2=1/Q$: every vector off that line
has error at least $1/Q$, and a Bezout pair with
$Qb+Pd=1$ realizes that error. Polynomial equivalence proves the
stated logarithmic limits. In the moment norm itself, two independent
integer pairs cannot both have zero actual output, so
\eqref{cal:eq:cal-existing} gives $\lambda_2\ge1/Q$.
\end{proof}

\paragraph{Arithmetic meaning.}
The theorem separates an experimentally small second minimum from a
proved arithmetic exponent. If the irrationality of $V$ were established,
its second-minimum lower decay exponent would determine $\mu(V)$ exactly.
In particular a rigorously proved value strictly below $1/2$ would imply
$\mu(V)>2$ and, by Roth's theorem~\cite{Roth55}, transcendence. The present paper proves
neither that inequality nor irrationality. The quantified lattice is an
exact reformulation of this additional Diophantine difficulty, not its
resolution.


% New consequences of the primitive coefficient certificates in Heisei R40 and Meiji R57.
% Needs amsmath, amssymb, amsthm and theorem/lemma/corollary environments.
\section{Algebraic root exclusion and the missing transcendence inequality}
\label{alg:sec:algebraic-exclusion}

The prime-scale determinant constructions have an arithmetic consequence
stronger than nonvanishing at a fixed rational argument. Their coefficient
certificates exclude all algebraic arguments of bounded absolute height,
uniformly in the degree. This does not by itself make their value at the
Gieseking constant small.

For a nonzero polynomial $P=\sum_{j=0}^s c_jX^j\in\mathbb Z[X]$, write
$H(P)=\max_j|c_j|$. For an algebraic number $\xi$ of degree $D$, with
primitive integral minimal polynomial
$A\prod_{i=1}^D(X-\xi_i)$, $A>0$, use the absolute logarithmic height
\[
 h(\xi)=\frac1D\left(\log A+
               \sum_{i=1}^D\log\max(1,|\xi_i|)\right).
\]
Every extension of the $p$-adic valuation below is normalized by $v_p(p)=1$.

\begin{theorem}[A uniform algebraic root-height certificate]
\label{alg:thm:root-height}
Let $p$ be a prime and $r\ge1$ an integer. Suppose
$P=\sum_{j=0}^s c_jX^j\in\mathbb Z[X]$, $s\ge1$, satisfies
\begin{equation}\label{alg:eq:weighted-primitive}
 v_p(c_0)=0,\qquad v_p(c_j)\ge rj\quad(1\le j\le s),
 \qquad v_p(c_s)=rs.
\end{equation}
Then every root of $P$ in $\overline{\mathbb Q}_p$ has valuation exactly
$-r$. If $\xi$ is any algebraic root of $P$ and $D=[\mathbb Q(\xi):\mathbb Q]$,
the leading and constant coefficients of its primitive minimal polynomial
satisfy
\[
 v_p(A)=rD,\qquad v_p(a_0)=0.
\]
In particular,
\begin{equation}\label{alg:eq:height-exclusion}
 h(\xi)\ge r\log p.
\end{equation}
Thus $P(\xi)\ne0$ simultaneously for every algebraic $\xi$ with
$h(\xi)<r\log p$, with no restriction on its degree.
\end{theorem}

\begin{proof}
Let $z$ be a $p$-adic root and write $t=v_p(z)$. The constant coefficient
is nonzero, so $z\ne0$. If $t>-r$, the constant summand in $P(z)$ has
valuation zero, while every other summand has positive valuation. It cannot
cancel. If $t<-r$, the leading summand has valuation $s(r+t)$, and every
summand of degree $j<s$ has strictly larger valuation, since
$v_p(c_jz^j)\ge j(r+t)>s(r+t)$. This is impossible as well. Therefore $t=-r$.

Let $M(X)=A\prod_{i=1}^D(X-\xi_i)=\sum a_kX^k$ be the primitive minimal
polynomial of an algebraic root. Since $M$ divides $P$ over $\mathbb Q$,
every $p$-adic conjugate $\xi_i$ is a root of $P$ and has valuation $-r$.
The constant coefficient has exact valuation
$v_p(a_0)=v_p(A)-rD$, while every elementary symmetric expression gives
\[
 v_p(a_k)\ge v_p(A)-r(D-k)=v_p(a_0)+rk.
\]
Integrality gives $v_p(a_0)\ge0$. If it were positive, every coefficient
would be divisible by $p$, contrary to primitivity. Hence $v_p(a_0)=0$
and $v_p(A)=rD$. Since $A\ge p^{rD}$, the defining height formula proves
\eqref{alg:eq:height-exclusion}.
\end{proof}

\begin{corollary}[Application to both prime-scale constructions]
\label{alg:cor:determinant-root-height}
The primitive polynomial $G_p$ of Heisei R40~\cite{H40}, of degree $(p-1)/2$, and
the primitive polynomial $G_{a,p}$ of Meiji R57, of degree $p$ at a prime
good for its fixed positive shift $a$, satisfy
\[
 G_p(\xi)\ne0,\qquad G_{a,p}(\xi)\ne0
 \quad\hbox{whenever}\quad h(\xi)<2\log p.
\]
In each case every algebraic root has height at least $2\log p$.
In particular, for every fixed algebraic $\xi$, either family is nonzero
at $\xi$ for every sufficiently large eligible prime.
\end{corollary}
\begin{proof}
Both source constructions prove
$v_p(c_0)=0$, $v_p(c_j)\ge2j$, and equality $v_p(c_s)=2s$ at the leading
coefficient. Apply Theorem~\ref{alg:thm:root-height} with $r=2$.
The good-prime restriction in Meiji R57 excludes only finitely many primes
for each fixed positive shift. The shift is fixed before $p$ varies.
\end{proof}

This is a consequence of the already established coefficient certificates,
not a new proof of those determinant constructions or a claim of a new
general height principle. It closes their algebraic-specialization
nonvanishing question, including all degrees and heights below the stated
threshold. An arbitrary rational rescaling before primitive normalization
does not preserve the coefficient hypothesis and must not be substituted.

\begin{lemma}[A norm lower bound with the degree explicit]
\label{alg:lem:algebraic-evaluation}
Let $\xi$ be algebraic of degree $D$ and let
$0\ne P\in\mathbb Z[X]$ have degree $s$, with $P(\xi)\ne0$. Then
\begin{equation}\label{alg:eq:resultant-bound}
 \log|P(\xi)|\ge -(D-1)\log H(P)
 -(D-1)\log(s+1)-Ds\,h(\xi).
\end{equation}
\end{lemma}
\begin{proof}
Set $\xi_1=\xi$ and let $M$ be the primitive minimal polynomial above.
Its nonzero resultant with $P$ is an integer, so
\[
 1\le |\operatorname{Res}(M,P)|
   =A^s\prod_{i=1}^D|P(\xi_i)|.
\]
For $i\ge2$,
$|P(\xi_i)|\le(s+1)H(P)\max(1,|\xi_i|)^s$.
Also
$A\prod_{i=2}^D\max(1,|\xi_i|)\le\exp(Dh(\xi))$.
Isolating the first factor and taking logarithms gives the assertion.
\end{proof}

\begin{theorem}[A degree-exclusion and transcendence criterion]
\label{alg:thm:degree-exclusion}
Let $x$ be a real number. For infinitely many primes $p\to\infty$, let
$P_p\in\mathbb Z[X]$ be nonzero polynomials which satisfy
\eqref{alg:eq:weighted-primitive} for some fixed $r\ge1$.
Suppose there is a positive scale $T_p\to\infty$ such that
\[
 \deg P_p=o(T_p),\qquad
 \log H(P_p)\le B T_p+o(T_p),\qquad
 \log|P_p(x)|\le-A T_p+o(T_p),
\]
where $A>0$, $B\ge0$, and $\log0=-\infty$.
Then $x$ cannot be algebraic of a degree $D$ for which
\begin{equation}\label{alg:eq:degree-threshold}
 (D-1)B<A.
\end{equation}
In particular, $B=0$ implies that $x$ is transcendental.
\end{theorem}
\begin{proof}
If $x$ were algebraic, Theorem~\ref{alg:thm:root-height} would make $P_p(x)$
nonzero for all sufficiently large primes in the sequence. Apply
Lemma~\ref{alg:lem:algebraic-evaluation}. Its last term is $o(T_p)$ because
$x$ and $D$ are fixed, and its logarithmic degree term is also $o(T_p)$.
The result is
$\log|P_p(x)|\ge-(D-1)BT_p-o(T_p)$, contradicting the assumed upper
bound whenever \eqref{alg:eq:degree-threshold} holds.
\end{proof}

For the two constructions in Corollary~\ref{alg:cor:determinant-root-height},
the degrees are $O(p)$, so $T_p=p^2$ is permissible. A hypothetical bound
$\log|P_p(V)|\le-Ap^2+o(p^2)$ would prove irrationality without a separate
height bound: take $D=1$ in \eqref{alg:eq:resultant-bound}. To exclude quadratic
values as well one needs, for example, $\log H(P_p)\le Bp^2+o(p^2)$ with
$A>B$. To establish transcendence by this route, it would suffice to prove
the same strict quadratic decay together with $\log H(P_p)=o(p^2)$.
None of these required decay estimates is established for the present
primitive determinant families. Their known local nonvanishing and formal
normality do not provide the missing inequality at the real place.

\paragraph{The external arithmetic result and the field of scalars.}
Calegari, Dimitrov and Tang prove the $\mathbb Q$-linear independence of
$1,\pi^2,L(2,\chi_{-3})$~\cite{CDT24}. Their later effective result gives the bound
$|L(2,\chi_{-3})-u/v|>v^{-24781}$ outside a finite computable set of reduced
rational approximants~\cite{CDT25}. Neither statement supplies
$\mathbb Q(\sqrt3)$-linear independence. In particular, if $V$ were a
nonzero rational number, then
$L(2,\chi_{-3})=4V/(3\sqrt3)$ would be a quadratic irrational, consistent
with both statements. Thus the classical irrationality of the scaled
constant cannot replace the real-place estimate in
Theorem~\ref{alg:thm:degree-exclusion}.



\section{What is closed, and what remains}
The factor-three inverse theorem closes an explicit denominator question.
The first-layer theorem closes an image classification in its no-wrap
range and, with the nested-image input, a missing prime interval.
The dyadic index determines an intrinsic arithmetic cost which cannot
be removed by changing the rational output basis. The positive-family
optimization is a global theorem for its stated budget, and the calibrated
lattice theorem identifies its exact Diophantine content.

These results also prevent several invalid shortcuts. Rational surjectivity
does not produce integral controls of comparable size. A positive rational
kernel does not bound the height of an integral representative. A nonzero
formal coordinate vector need not give a nonzero real form. A small
unnormalized determinant can become large after primitive normalization.
A bound for one successive minimum cannot replace a bound for several
independent directions.

Two concrete infinite problems remain sufficient for an arithmetic conclusion.
One is the existence of the positive, nuisance-cancelling, integral-output
controls of Corollary~\ref{opt:taisho:newcriterion} within its proved height
budget. The other is strict quadratic real-place decay of a primitive
prime-scale polynomial, with the coefficient height needed for the desired
degree exclusion in Theorem~\ref{alg:thm:degree-exclusion}. These are genuine
additional estimates, not consequences of the local algebra proved here.

There is a useful distinction between the two basis results. The inverse
denominator theorem concerns rational coefficients in the complete $4n$
window and improves their arithmetic cost. The dyadic obstruction concerns
a different odd space and a complete collection of integer-output controls.
They cannot be identified in either direction. Similarly the new optimized
parameter does not inherit a positive-cone theorem proved at a different
parameter without another argument.

\appendix
\section{Source coverage and reproducibility}
The supplied corpus contains the following Gieseking manuscripts. The seven
Mills/Pisot notes in the additional archive are outside the present synthesis.
\begin{center}
\begin{tabular}{lrrr}\toprule
Series & Earlier notes & Additional notes & Total\\\midrule
Keio &22&8&30\\
Reiwa &20&12&32\\
Meiji &37&9&46\\
Taisho &24&6&30\\
Heisei &33&12&45\\\midrule
Total &136&47&183\\\bottomrule
\end{tabular}
\end{center}
The companion verification archive includes a title/path/hash catalogue,
five series audits, the arithmetic audit, and exact checkers. The audits
state which proof mechanisms were examined closely and which older results
were superseded. This is not a claim that every historical calculation or
every source-mentioned external checker was independently rerun.

The finite replays include: the complete mixed-moment determinant through
$n=5$; first and second prime-degree thresholds and simple/double pole
witnesses for 82 pairs $(p,m)$; 95 complete saturation-source spaces;
the dyadic index through $n=25$ and the full odd output identities through
$n=10$; conic coordinate denominator and normalizer checks; and exact rational
enclosures for the optimization constants. The new normalized inverse
estimates are also checked independently on finite moment matrices.
No finite replay is used as a substitute for an all-index proof.

For the new theorems, separate proof reviews covered the factor-three carry
argument, the dyadic Smith index, the saturation interval, the positive
budget optimization, and the calibrated minima. These are internal checks,
not external peer review or a certification of historical novelty.
The published results cited below are used only with their stated scalar
fields and hypotheses.

\begin{thebibliography}{99}\small\raggedright
\bibitem{CDT24} F. Calegari, V. Dimitrov and Y. Tang,
\emph{The linear independence of $1$, $\zeta(2)$, and $L(2,\chi_{-3})$},
arXiv:2408.15403v2 (2024). \url{https://arxiv.org/abs/2408.15403}.
\bibitem{CDT25} F. Calegari, V. Dimitrov and Y. Tang,
\emph{Arithmetic holonomy bounds and effective Diophantine approximation},
arXiv:2510.04156 (2025), 2026 SIAM proceedings version; Theorem 5.1.
\url{https://arxiv.org/abs/2510.04156}.
\bibitem{Roth55} K. F. Roth, \emph{Rational approximations to algebraic numbers},
Mathematika \textbf{2} (1955), 1--20.
\url{https://doi.org/10.1112/S0025579300000644}.
\bibitem{Corpus} Keio, Reiwa, Meiji, Taisho and Heisei series,
\emph{Supplied Gieseking research corpus}: 136 earlier notes and 47 subsequent
notes (2026). Unpublished manuscripts supplied for this synthesis; individual
identifiers below are corpus labels, not journal citations.
\bibitem{T05} Taisho R05, \emph{A Five-Constant Moment Algebra for the Gieseking Volume}.
\bibitem{T13} Taisho R13, Frobenius moment descent, supplied manuscript.
\bibitem{T28} Taisho R28, low-degree local images and boundary model, supplied manuscript.
\bibitem{T36} Taisho R36, \emph{Arithmetic and Decay Bounds for Proportional Bernstein--Rodrigues Controls}.
\bibitem{T37} Taisho R37, \emph{Uniform Positive Bernstein Cones and a Three-Dimensional Weighted Lattice Reduction}.
\bibitem{T38} Taisho R38, \emph{Sharp Prime-Degree Thresholds for Mixed Gieseking Period Images}.
\bibitem{M58} Meiji R58, \emph{An Exact Four-Block Determinant and Minimal-Window Gieseking Moment Normality}.
\bibitem{M59} Meiji R59, \emph{A Calibrated Two-Dimensional Moment Lattice}.
\bibitem{M61} Meiji R61, \emph{An Explicit Full-Window Moment Inverse with Exponential-Height Rational Coefficients}.
\bibitem{M62} Meiji R62, \emph{Laurent and Centered-Angle Primitives for the Full Gieseking Moment Inverse}.
\bibitem{M57} Meiji R57, \emph{Positive Block Shifts Remove a Prime-Scale Gieseking Determinant Rank Obstruction}.
\bibitem{H06} Heisei R06, \emph{An Exact Determinant and All-Size Normality for Gieseking Moment Cancellation}.
\bibitem{H40} Heisei R40, primitive prime-scale determinant construction, supplied manuscript.
\bibitem{H42} Heisei R42, continuous-row/common-positive-weight odd-column sign regularity, supplied manuscript.
\bibitem{H43} Heisei R43, \emph{An Explicit Binomial Adjoint Factorization and Small Denominators in Tail Gieseking Rows}.
\bibitem{H48} Heisei R48v2, direct odd rational-output construction and coefficient denominators, supplied manuscript.
\bibitem{H49} Heisei R49, \emph{The Sharp Large-Prime Residue Rank of the Odd Rational-Output Basis}.
\bibitem{H50} Heisei R50, \emph{The Exact Asymptotic Denominator of the Direct Odd Output Basis}.
\bibitem{K49} Keio R49, \emph{Normality of Three Consecutive Legendre Density Periods}.
\bibitem{K50} Keio R50, \emph{Four-Consecutive Conic Normality and All-Index High-Block Normalizers}.
\bibitem{K51} Keio R51, \emph{A New Endpoint Source and a Two-Block Local Image}.
\bibitem{R49} Reiwa R49, complete coefficient-saturation basis and Frobenius descent, supplied manuscript.
\bibitem{R52} Reiwa R52, two-mode saturation bands and the remaining two-block interval, supplied manuscript.
\end{thebibliography}
\end{document}
